% content/c03-hilbertraum/en.tex — Chapter 3: The Hilbert-space toolkit
% (English edition; the German de.tex leads, same constructs, IDs and displays)
%
% Scope: the material of the second, third and fourth lectures (latex/lectures/de/vl02.tex,
% vl03.tex and vl04.tex, with the vl03 subsection on operators). New sections go between the
% last core section and the Fundament blocks at the end.
% Maps: evidence/migrations/maps/vl02-map.json, vl03-map.json, vl04-map.json (lane E). Defects:
% M-0013 to M-0016 (vl02), M-0017 to M-0024 (vl03), M-0005 and docs/plan/lec34/NOTES-author.md
% (vl04). Definitions D75 to D80, D81 to D86 and D14, D87 to D94, D17, D18, D20, convention
% CONV-HAT; derivations evidence/derivations/c03/DERIVATION.md; Reed-Simon page-image rows P-0140
% to P-0150 (review receipt R-0140), P-0152 to P-0164 (review receipt R-0152), from P-0318 (vl04,
% review receipt pending). Reviews of Lecture 3: docs/review/lectures/REVIEW-lec3.md,
% evidence/derivations/c03/REVIEW-lec3.md and REVIEW-lec3-second.md. G-0072 (a) and (b) applied
% after TJO's answer of 2026-10-29. Hats on every operator (D-HAT), except the identity, states and
% the Pauli matrices (CONV-HAT; none of the latter two occurs in c03).

\chapter{The Hilbert-space toolkit}\label{c03}

To describe the states of quantum mechanics, the \emph{Hilbert space} will play a major role.
It is the mathematical core of quantum mechanics: without it we cannot even talk about the
objects of the theory. To define this space we first need \emph{complex vector spaces}. Such
spaces are defined in complete analogy with the real case, which you know from Theoretical
Physics~A \srcref{tpc-script-v0.2}{vl02, introduction}.

Why \emph{complex} and not real vectors appear here, we do not know: nature is simply built
that way, we have to live with it, and no better alternative has been found so far. This fact
has very important physical consequences, which you will see in the course of these lectures.

The qubit has already appeared in Chapter~\ref{c02}, with its ball of states. How exactly the
states of quantum mechanics are built from a Hilbert space is the subject of Chapter~4. This
chapter provides the tools, in a form you can hold in your hand: columns of complex numbers.

\section{Complex vectors}\label{c03.s1}

Fortunately the definition is simple. You know complex numbers, and you know column vectors;
now the two are brought together.

\tpcdef{D75}

A complex vector is therefore nothing but a list of complex numbers. On the blackboard I write
it in bold and underlined; in print it is just bold, $\vb{v}$. If you multiply a complex
vector by a scalar or add two of them, you get a complex vector again.

For the vector space itself these two rules suffice. You do not need an inner product for it:
the dagger operation and the inner product of the next section are extra structure, and only
with them does $\mathbb{C}^n$ become a Hilbert space.

The simplest case, $\mathbb{C}^1$, is boring: these are just the complex numbers. The first
interesting example is $\mathbb{C}^2$. Computationally this space is very simple, and
physically it is nevertheless central: two-dimensional complex vectors describe the states of
a qubit, the basic building block of a quantum computer, and one needs them to describe the
spin of an electron \srcref{tpc-script-v0.2}{vl17, section on spin}. More precisely, a
vector describes a \emph{pure} state; mixed states, such as the points inside the ball of
states of Chapter~\ref{c02}, need density matrices, and these come in Chapter~4
\srcref{tpc-script-v0.2}{vl06, Postulate 2}.

\begin{example}[Computing in $\mathbb{C}^2$]\label{c03.ex.c2}\claim{c03.cl.c2}{derived}
  The vector space $\mathbb{C}^2$ consists of all vectors with two complex entries. Let
  \begin{equation}\label{c03.eq.c2-vektoren}
    \vb{v} = \begin{pmatrix} 1+i \\ 1-i \end{pmatrix},
    \qquad
    \vb{w} = \begin{pmatrix} 1 \\ -1 \end{pmatrix}
  \end{equation}
  be two vectors; the first one deliberately has complex entries, to stress that this is
  allowed. Computing entry by entry,
  \begin{equation}\label{c03.eq.c2-rechnung}
    \vb{v} + \vb{w} = \begin{pmatrix} 2+i \\ -i \end{pmatrix},
    \qquad
    i\,\vb{v} = \begin{pmatrix} i + i^2 \\ i - i^2 \end{pmatrix}
              = \begin{pmatrix} -1+i \\ 1+i \end{pmatrix} .
  \end{equation}
\end{example}

\section{The dagger operation and the inner product}\label{c03.s2}

A complex vector space on its own carries no further structure. To get to the Hilbert space we
need additional operations. The first is the \emph{dagger operation}: a column vector becomes
a row vector, and the entries are complex conjugated on the way. With it the inner product can
be introduced at once: apply the dagger operation to the first vector and multiply the row by
the column of the second. The result is always a complex number.

\tpcdef{D76}

\begin{tpcfigure}{c03.fig.kreuz}
  \demo{c03.demo.kreuz}
  \caption{The dagger operation and the inner product. (a) The entries of a vector in
    $\mathbb{C}^2$ as points of the complex plane, here $v_1 = 2+i$ and $v_2 = -1+2i$ (filled).
    The dagger operation reflects every entry in the real axis (open circles,
    $v_1^* = 2-i$ and $v_2^* = -1-2i$) and lays the column down as a row. (b) The inner
    product of Example~\ref{c03.ex.skalarprodukt}: the row $\vb{v}\adj$ is multiplied entry by
    entry with the column $\vb{w}$, and the products are added.}
\end{tpcfigure}

You do not have to do the complex conjugation in the dagger operation; you could simply lay
a column down. It pays off, however, as soon as you multiply a vector with itself. Since
$v_j^* v_j = \abs{v_j}^2 = (\operatorname{Re} v_j)^2 + (\operatorname{Im} v_j)^2$, we have
\begin{equation}\label{c03.eq.selbst}\claim{c03.cl.selbst}{derived}
  \vb{v}\adj\vb{v} = \sum_{j=1}^{n} v_j^* v_j = \sum_{j=1}^{n} \abs{v_j}^2 \;\geq\; 0 .
\end{equation}
The inner product of a vector with itself is therefore a real number and never negative.

\begin{example}[Without conjugation]\label{c03.ex.ohne-stern}\claim{c03.cl.ohne-stern}{derived}
  Let $\vb{v} = \begin{pmatrix} 1 \\ i \end{pmatrix}$. Without conjugation, row times column
  would give $1 \cdot 1 + i \cdot i = 1 + i^2 = 0$, although $\vb{v}$ is not the zero vector.
  With the dagger operation,
  \begin{equation}\label{c03.eq.ohne-stern}
    \vb{v}\adj\vb{v} = 1 \cdot 1 + (-i) \cdot i = 1 + 1 = 2 .
  \end{equation}
  Without the star, a vector that is not zero would have “length” zero.
\end{example}

The inner product is the extra structure with which we shall later compute
\emph{probabilities} in quantum mechanics. Quantum mechanics is a probabilistic theory
(Chapter~\ref{c01}) and therefore needs a method to compute probabilities; the inner product
will be that method \srcref{tpc-script-v0.2}{vl03, section on inner products}.

\begin{example}[An inner product in $\mathbb{C}^2$]\label{c03.ex.skalarprodukt}\claim{c03.cl.vw}{derived}
  For the vectors of Equation~\eqref{c03.eq.c2-vektoren} one computes in two steps, first the
  dagger operation, then row times column:
  \begin{equation}\label{c03.eq.vw}
    \vb{v}\adj\vb{w}
    = \begin{pmatrix} 1-i & 1+i \end{pmatrix} \begin{pmatrix} 1 \\ -1 \end{pmatrix}
    = (1-i) - (1+i) = -2i .
  \end{equation}
  Does it matter which of the two vectors gets the dagger? Yes: the \emph{first} vector
  always gets the dagger. In the other order the conjugation is trivial here, because
  $\vb{w}$ is real, and the result is different:
  \begin{equation}\label{c03.eq.wv}
    \vb{w}\adj\vb{v}
    = \begin{pmatrix} 1 & -1 \end{pmatrix} \begin{pmatrix} 1+i \\ 1-i \end{pmatrix}
    = (1+i) - (1-i) = 2i .
  \end{equation}
  The two results are complex conjugates of each other. So the order matters, but the
  difference is under control.
\end{example}

That this is no accident can be read off from the sum. The following rules hold for every
$n$.

\begin{proposition}[Rules for the inner product in $\mathbb{C}^n$]\label{c03.prop.skalarprodukt}\claim[proof-included]{c03.cl.skalarprodukt}{derived}
  Let $\vb{v}, \vb{w}, \vb{u} \in \mathbb{C}^n$ and $\alpha, \beta \in \mathbb{C}$. Then:
  \begin{enumerate}
    \item Swapping conjugates: $\vb{w}\adj\vb{v} = (\vb{v}\adj\vb{w})^*$.
    \item Linear in the second factor:
      $\vb{v}\adj(\alpha\vb{w} + \beta\vb{u}) = \alpha\, \vb{v}\adj\vb{w} + \beta\, \vb{v}\adj\vb{u}$.
    \item Antilinear in the first factor:
      $(\alpha\vb{v} + \beta\vb{u})\adj\vb{w} = \alpha^*\, \vb{v}\adj\vb{w} + \beta^*\, \vb{u}\adj\vb{w}$.
    \item Positivity: $\vb{v}\adj\vb{v} \geq 0$, and $\vb{v}\adj\vb{v} = 0$ holds only for
      $\vb{v} = 0$.
  \end{enumerate}
\end{proposition}
\begin{proof}
  On~1: the conjugate of a sum is the sum of the conjugates, and $(z^*)^* = z$, so
  \begin{equation}\label{c03.eq.vertauschen}
    (\vb{v}\adj\vb{w})^* = \sum_{j=1}^{n} (v_j^* w_j)^* = \sum_{j=1}^{n} v_j\, w_j^*
    = \sum_{j=1}^{n} w_j^* v_j = \vb{w}\adj\vb{v} .
  \end{equation}
  On~2: $\sum_j v_j^* (\alpha w_j + \beta u_j) = \alpha \sum_j v_j^* w_j + \beta \sum_j v_j^* u_j$.
  On~3: the $j$-th entry of $(\alpha\vb{v} + \beta\vb{u})\adj$ is
  $(\alpha v_j + \beta u_j)^* = \alpha^* v_j^* + \beta^* u_j^*$; then as in~2. On~4:
  Equation~\eqref{c03.eq.selbst}. A sum of numbers $\abs{v_j}^2 \geq 0$ is zero only if every
  term is zero, that is, if every $v_j = 0$.
\end{proof}

When a factor is pulled out of the first vector, a star appears. That is the convention of
this course; some books do it the other way round and call the inner product linear in the
first vector \srcref{reed-simon-1}{p.~37}.

\begin{exercise}[Computing with the dagger]\label{c03.exc.rechnen}\claim{c03.cl.rechnen}{derived}
  Let $\vb{v} = \begin{pmatrix} i \\ 2 \end{pmatrix}$ and
  $\vb{w} = \begin{pmatrix} 1 \\ i \end{pmatrix}$. Compute $3\vb{v}$, $\vb{v} + \vb{w}$,
  $\vb{v}\adj\vb{w}$, $\vb{w}\adj\vb{v}$ and $\vb{v}\adj\vb{v}$, and check rule~1 of
  Proposition~\ref{c03.prop.skalarprodukt}.
\end{exercise}
\begin{solution}
  We have $3\vb{v} = \begin{pmatrix} 3i \\ 6 \end{pmatrix}$ and
  $\vb{v} + \vb{w} = \begin{pmatrix} 1+i \\ 2+i \end{pmatrix}$. With
  $\vb{v}\adj = \begin{pmatrix} -i & 2 \end{pmatrix}$ and
  $\vb{w}\adj = \begin{pmatrix} 1 & -i \end{pmatrix}$,
  \begin{equation}\label{c03.eq.rechnen}
    \vb{v}\adj\vb{w} = -i + 2i = i,
    \qquad
    \vb{w}\adj\vb{v} = i - 2i = -i,
    \qquad
    \vb{v}\adj\vb{v} = (-i)\,i + 4 = 5 .
  \end{equation}
  So $\vb{w}\adj\vb{v} = -i = (\vb{v}\adj\vb{w})^*$, as rule~1 says.
\end{solution}

\section{Infinitely many entries}\label{c03.s3}

At this point one could already define finite-dimensional Hilbert spaces. I deliberately
postpone that, however, because infinite-dimensional vector spaces play a very important role
in physics. The reason is a physical one: our best models are fields, and fields have
infinitely many degrees of freedom, as you know from Theoretical Physics~B
\srcref{tpc-script-v0.2}{vl02, paragraph on infinite-dimensional vector spaces}.

The core of this course nevertheless stays finite-dimensional up to and including Chapter~4:
everything said there about states and measurements takes place in $\mathbb{C}^n$. The
infinite-dimensional spaces enter the core only in Chapter~5, with the particle on a line.
If you want to see now what goes wrong when one simply sets $n = \infty$ in the formulas,
and how to repair it, you will find it in two deeper blocks. The first is about columns
with infinitely many entries\zoomref{c03.v1}.

\begin{vertiefung}[id=c03.v1,
    titel={Columns with infinitely many entries},
    wozu={If one simply sets $n = \infty$ in the formulas of Section~\ref{c03.s2}, the inner
      product can diverge. This block shows how to repair it: one admits only vectors whose
      inner product with themselves is finite. These vectors form the space $\elltwo$.},
    voraussetzung={Section~\ref{c03.s2}; convergent series. The core up to Chapter~4 does not
      need this block.},
    zurueck={A vector with infinitely many entries is admitted if the sum of the squared
      moduli of its entries is finite. For such vectors the inner product is defined as an
      infinite sum and obeys the same rules as in $\mathbb{C}^n$.}]
  One can picture infinite-dimensional vectors as columns with infinitely many entries, that
  is, as infinite lists of complex numbers:
  \begin{equation*}
    \vb{v} = \begin{pmatrix} v_1 \\ v_2 \\ v_3 \\ \vdots \end{pmatrix} .
  \end{equation*}
  Addition and multiplication by scalars go entry by entry, as in $\mathbb{C}^n$. For the
  inner product the obvious generalisation is the series
  $\vb{v}\adj\vb{w} = \sum_{j=1}^{\infty} v_j^* w_j$: as a rule, one sets $n = \infty$ in all
  equations. But here problems can arise, convergence problems. One shows up straight away.

  \begin{example}[The vector made up entirely of ones]\label{c03.ex.einsen}\claim{c03.cl.einsen}{derived}
    The vector made up entirely of ones, $v_j = 1$ for all $j$, is a perfectly legal
    infinite list of complex numbers. But its inner product with itself diverges: the
    partial sums are
    \begin{equation}\label{c03.eq.einsen}
      \sum_{j=1}^{N} \abs{v_j}^2 = \sum_{j=1}^{N} 1 = N \;\longrightarrow\; \infty
      \qquad (N \to \infty) .
    \end{equation}
  \end{example}

  If we want to describe the states of physical systems by such vectors, that is a problem,
  at the latest when it comes to probabilities. The answer is as simple as can be: we do not
  admit such vectors. In the mathematical literature this is described in a very complicated
  way; physically it is simple.

  \tpcdef{D77}

  We shall hardly need the name $\elltwo$ later, but you should have seen it once: the $2$
  comes from the square in the sum, and the sum runs over all natural numbers. Reed and Simon
  write the space with a subscript $2$ \srcref{reed-simon-1}{p.~40}. That the admitted
  vectors again form a vector space is not obvious, but true:

  \begin{proposition}[$\elltwo$ is a vector space with an inner product]\label{c03.prop.elltwo}\claim[proof-included]{c03.cl.elltwo}{derived}
    Let $\vb{v}, \vb{w} \in \elltwo$ and $a \in \mathbb{C}$. Then:
    \begin{enumerate}
      \item $a\vb{v} \in \elltwo$ and $\vb{v} + \vb{w} \in \elltwo$.
      \item The series $\vb{v}\adj\vb{w} = \sum_{j=1}^{\infty} v_j^* w_j$ converges absolutely,
        and $\sum_{j=1}^{\infty} \abs{v_j^* w_j} \leq \tfrac12\, \vb{v}\adj\vb{v} + \tfrac12\, \vb{w}\adj\vb{w}$.
      \item The four rules of Proposition~\ref{c03.prop.skalarprodukt} hold in $\elltwo$.
    \end{enumerate}
  \end{proposition}
  \begin{proof}
    For complex numbers $a, b$ we have $0 \leq (\abs{a} - \abs{b})^2$ and
    $\abs{a + b}^2 = \abs{a}^2 + \abs{b}^2 + 2 \operatorname{Re}(a^* b)$ with
    $\operatorname{Re}(a^* b) \leq \abs{a}\,\abs{b}$. From these follow
    \begin{equation}\label{c03.eq.ab}
      \abs{a^* b} = \abs{a}\,\abs{b} \leq \tfrac12 \abs{a}^2 + \tfrac12 \abs{b}^2,
      \qquad
      \abs{a + b}^2 \leq 2\abs{a}^2 + 2\abs{b}^2 .
    \end{equation}
    On~1: $\sum_{j} \abs{a v_j}^2 = \abs{a}^2 \sum_j \abs{v_j}^2$. By the second inequality
    every partial sum $\sum_{j \leq N} \abs{v_j + w_j}^2$ is at most
    $2\, \vb{v}\adj\vb{v} + 2\, \vb{w}\adj\vb{w}$. The partial sums increase with $N$ and are
    bounded, so they converge. On~2: in the same way with the first inequality; an absolutely
    convergent series of complex numbers converges. On~3: the rules hold for every partial
    sum up to $N$, since that is a computation in $\mathbb{C}^N$. Limits are compatible with
    sums, multiples and conjugation. If $\vb{v}\adj\vb{v} = 0$, every partial sum
    $\sum_{j \leq N} \abs{v_j}^2$ is zero, hence every $v_j = 0$.
  \end{proof}

  With this condition nothing can become infinitely large any more when probabilities are
  computed. A sharper bound than the one in part~2, the Cauchy-Schwarz inequality, is in the
  Foundations\zoomref{c03.f1}.

  By the way, it is not enough for the entries to go to zero. What matters is how fast.

  \begin{example}[The rate is what matters]\label{c03.ex.rate}\claim{c03.cl.rate}{derived}
    The vector with $v_j = 1/j$ lies in $\elltwo$, the vector with $u_j = 1/\sqrt{j}$ does
    not, although its entries go to zero as well. Indeed, for every $N \geq 2$,
    \begin{equation}\label{c03.eq.rate}
      \begin{aligned}
        \sum_{j=1}^{N} \frac{1}{j^2}
          &\leq 1 + \sum_{j=2}^{N} \Bigl( \frac{1}{j-1} - \frac{1}{j} \Bigr)
           = 2 - \frac{1}{N} < 2, \\
        \sum_{j=1}^{N} \abs{u_j}^2 = \sum_{j=1}^{N} \frac{1}{j}
          &\geq \int_{1}^{N+1} \frac{\dd x}{x} = \ln(N+1) \;\longrightarrow\; \infty .
      \end{aligned}
    \end{equation}
    In the first line $1/j^2 \leq 1/((j-1)j) = 1/(j-1) - 1/j$, and the sum on the right
    telescopes. In the second, the term $1/j$ lies above the area under $1/x$ between $j$ and
    $j+1$, because $1/x$ decreases. The partial sums of the squared moduli of $\vb{u}$ thus
    grow beyond all bounds, at least like a logarithm.
  \end{example}

  \begin{tpcfigure}{c03.fig.partialsummen}
    \demo{c03.demo.partialsummen}
    \caption{The admission test for $\elltwo$: the partial sums
      $\sum_{j=1}^{N} \abs{v_j}^2$ for $N = 1, \dots, 25$, rounded to four places; the curve
      for the vector of all ones is drawn only up to $N = 5$. For $v_j = 1/j$
      (dots, solid) they stay below the bound $2$ of Equation~\eqref{c03.eq.rate} (thin
      line): the vector is admitted. For $v_j = 1/\sqrt{j}$ (dots, dashed) they grow beyond
      all bounds, at least like $\ln(N+1)$, for the vector made up entirely of ones (dots, dash-dotted) like $N$: neither
      is admitted, although in the middle case the entries go to zero. In $\mathbb{C}^n$ the
      sum has only $n$ terms, and there is nothing to test.}
  \end{tpcfigure}

  This space is not just mathematics. $\elltwo$ is the state space of the harmonic
  oscillator, which we look at closely in Chapter~7. As a quantum physicist I work
  practically every day with these two vector spaces, with $\mathbb{C}^2$ and with
  $\elltwo$.
\end{vertiefung}

The third space we need consists not of columns but of functions: to every point $x$ of the
real axis a complex number $f(x)$ is assigned. Such functions are very similar to the fields
of Theoretical Physics~B, the electric field or the potential, except that the values are now
complex. They describe the pure states of a spinless particle in one
dimension\zoomref{c03.v2}.

\begin{vertiefung}[id=c03.v2,
    titel={Functions as vectors},
    wozu={Functions can be added and multiplied by numbers like vectors. Here they get an
      inner product, with an integral in place of the sum (with one reservation, which the
      Foundations resolve), and the same admission condition as the columns in
      $\elltwo$.},
    voraussetzung={Section~\ref{c03.s2}; improper integrals. Block~\ref{c03.v1} helps but is
      not needed. The core up to Chapter~4 does not need this block; Chapter~5 introduces the
      space in the core once more.},
    zurueck={Square-integrable functions form a vector space on which the integral
      $\braket{f}{g} = \int f^*(x)\, g(x)\, \dd x$ is defined; it becomes a true inner
      product only once functions that agree almost everywhere are identified. The constant function
      $1$ does not belong to it.}]
  A function $f\colon \mathbb{R} \to \mathbb{C}$ assigns a complex number to every point
  $x$. That is somewhat hard to draw: a complex number needs two axes, so the function needs
  a third dimension (Figure~\ref{c03.fig.komplex}). Functions are added point by point,
  $(f + g)(x) = f(x) + g(x)$, and likewise $(af)(x) = a f(x)$ for a scalar $a$. In this sense
  a function is a vector with one entry $f(x)$ for every point $x$. The integral over $x$
  takes the place of the sum over $j$.

  \begin{tpcfigure}{c03.fig.komplex}
    \demo{c03.demo.komplex}
    \caption{A complex-valued function needs three dimensions. Drawn is
      $f(x) = e^{-x^2/2}\, e^{2ix}$ for $-3 \leq x \leq 3$: to every $x$ belongs the point
      $(\operatorname{Re} f(x), \operatorname{Im} f(x))$ of a complex plane standing
      perpendicular to the $x$-axis. The curve winds around the $x$-axis; its distance from
      the axis is $\abs{f(x)} = e^{-x^2/2}$ (thin circles at $x = -1, 0, 1$; thin spokes join the axis to the curve). Dashed is the real part alone,
      in the plane $\operatorname{Im} f = 0$: the usual drawing, which lacks half of the
      information.}
  \end{tpcfigure}

  In this block all functions are piecewise continuous, so that every integral is an ordinary
  improper integral. What changes when this is dropped is in the Foundations of this
  chapter.

  \tpcdef{D78}

  The sign $\braket{f}{g}$ stands for a number here, just as $\vb{v}\adj\vb{w}$ does in
  $\mathbb{C}^n$. This course writes every inner product this way; what the vertical bar
  means becomes clear with Dirac's notation.

  \begin{proposition}[Square-integrable functions form a vector space]\label{c03.prop.ltwo}\claim[proof-included]{c03.cl.ltwo}{derived}
    Let $f$ and $g$ be square-integrable and $a \in \mathbb{C}$. Then $af$ and $f + g$ are
    square-integrable, and the integral
    $\braket{f}{g} = \int_{-\infty}^{\infty} f^*(x)\, g(x)\, \dd x$ converges absolutely, with
    \begin{equation}\label{c03.eq.ltwo}
      \int_{-\infty}^{\infty} \abs{f^*(x)\, g(x)}\, \dd x
      \leq \tfrac12 \braket{f}{f} + \tfrac12 \braket{g}{g} .
    \end{equation}
    Rules~1 to~3 of Proposition~\ref{c03.prop.skalarprodukt} hold correspondingly, and
    $\braket{f}{f} \geq 0$.
  \end{proposition}
  \begin{proof}
    For complex numbers $a, b$ we have $\abs{a^* b} \leq \tfrac12 \abs{a}^2 + \tfrac12 \abs{b}^2$,
    because $0 \leq (\abs{a} - \abs{b})^2$, and $\abs{a + b}^2 \leq 2\abs{a}^2 + 2\abs{b}^2$,
    because $\abs{a + b}^2 = \abs{a}^2 + \abs{b}^2 + 2 \operatorname{Re}(a^* b)$ and
    $\operatorname{Re}(a^* b) \leq \abs{a}\,\abs{b}$. Insert $a = f(x)$ and $b = g(x)$; the
    inequalities then hold at every point $x$. Integrated over $[-X, X]$ they give bounds that do not depend on $X$.
    Integrals of non-negative functions over $[-X, X]$ increase with $X$ and are bounded, so
    they converge as $X \to \infty$. An absolutely convergent improper integral converges.
    Rules~1 to~3 hold for the integrand at every point, and the integral is linear. Finally,
    $\abs{f(x)}^2 \geq 0$ at every point.
  \end{proof}

  Rule~4 of Proposition~\ref{c03.prop.skalarprodukt} is missing from this list on purpose.
  That $\braket{f}{f} = 0$ holds only for $f = 0$ is not quite right for functions, and the
  repair is in the Foundations\zoomref{c03.f2}.

  \begin{example}[A non-example]\label{c03.ex.eins}\claim{c03.cl.eins}{derived}
    The constant function $f(x) = 1$ is not square-integrable:
    \begin{equation}\label{c03.eq.eins}
      \int_{-X}^{X} \abs{f(x)}^2\, \dd x = \int_{-X}^{X} 1\, \dd x = 2X
      \;\longrightarrow\; \infty
      \qquad (X \to \infty) .
    \end{equation}
    The area under the curve is infinite. This is exactly the counterpart of the vector made
    up entirely of ones in Example~\ref{c03.ex.einsen}.
  \end{example}

  \begin{example}[Going to zero is not enough]\label{c03.ex.gauss}\claim{c03.cl.gauss}{derived}
    The function $g(x) = e^{-x^2/2}$ is square-integrable, and so is the function
    $f(x) = e^{-x^2/2}\, e^{2ix}$ of Figure~\ref{c03.fig.komplex}, because
    $\abs{e^{2ix}} = 1$ and hence $\abs{f(x)}^2 = \abs{g(x)}^2 = e^{-x^2}$. With the
    substitution $x = y/\sqrt{2}$ and the Gaussian integral of Chapter~\ref{c01},
    $\int_{-\infty}^{\infty} e^{-y^2/2}\, \dd y = \sqrt{2\pi}$,
    \begin{equation}\label{c03.eq.gauss}
      \braket{g}{g} = \int_{-\infty}^{\infty} e^{-x^2}\, \dd x
      = \frac{1}{\sqrt{2}} \int_{-\infty}^{\infty} e^{-y^2/2}\, \dd y
      = \sqrt{\pi} .
    \end{equation}
    The function $h(x) = 1/\sqrt{1 + \abs{x}}$, by contrast, goes to zero as
    $\abs{x} \to \infty$ and is nevertheless not square-integrable:
    \begin{equation}\label{c03.eq.log}
      \int_{-X}^{X} \frac{\dd x}{1 + \abs{x}} = 2 \int_{0}^{X} \frac{\dd x}{1 + x}
      = 2 \ln(1 + X) \;\longrightarrow\; \infty .
    \end{equation}
    These examples show: that a function goes to zero is not enough. What decides is only
    whether the integral of its squared modulus is finite.
  \end{example}

  \begin{tpcfigure}{c03.fig.partialintegrale}
    \demo{c03.demo.partialintegrale}
    \caption{The same test for functions: the integral $\int_{-X}^{X} \abs{f(x)}^2\, \dd x$
      as a function of $X$. For $f(x) = e^{-x^2/2}$ (solid) it approaches $\sqrt{\pi}$ (thin
      line, Equation~\eqref{c03.eq.gauss}); for $f(x) = 1/\sqrt{1 + \abs{x}}$ (dashed) it
      grows like $2\ln(1 + X)$, for $f = 1$ (dash-dotted) like $2X$
      (Equations~\eqref{c03.eq.log} and~\eqref{c03.eq.eins}). The lines show continuously
      what Figure~\ref{c03.fig.partialsummen} shows for the discrete index $j$.}
  \end{tpcfigure}
\end{vertiefung}

Up to and including Chapter~4 the core works with finite-dimensional spaces $\mathbb{C}^n$;
for the qubit this is $\mathbb{C}^2$. In the lecture I put three spaces and three physical
systems side by side at this point, among them the two infinite-dimensional spaces of the
deeper blocks\zoomref{c03.v3}.

\begin{vertiefung}[id=c03.v3,
    titel={Three spaces, three systems},
    wozu={The three spaces of this chapter side by side, each with the physical system it
      belongs to. With them we have practically all the Hilbert spaces of the course.},
    voraussetzung={Blocks~\ref{c03.v1} and~\ref{c03.v2}. The core up to Chapter~4 does not
      need this block.},
    zurueck={Up to and including Chapter~4 the core works with $\mathbb{C}^n$; for the qubit
      this is $\mathbb{C}^2$.}]
  \begin{greybox}[Three spaces, three systems]
  With $\mathbb{C}^2$ and two infinite-dimensional spaces we have practically all the Hilbert
  spaces we shall need in this course: $\elltwo$, the columns with infinitely many entries
  whose squared moduli have a finite sum, and $\Ltwo{\mathbb{R}}$, formed from the functions
  on the line whose squared modulus has a finite integral, where functions that agree almost
  everywhere count as the same vector (Foundations of this chapter). Each space belongs to a
  physical system:
  \begin{equation}\label{c03.eq.drei-raeume}\claim{c03.cl.drei-raeume}{source-backed}
    \begin{aligned}
      &\mathbb{C}^2 && \text{\tpcifenglish{qubit}{Qubit}} \\
      &\elltwo && \text{\tpcifenglish{harmonic oscillator}{harmonischer Oszillator}} \\
      &\Ltwo{\mathbb{R}} && \text{\tpcifenglish{spinless particle in one dimension}{Teilchen ohne Spin in einer Dimension}}
    \end{aligned}
  \end{equation}
  The first space is finite-dimensional. The other two enter the core from Chapter~5 on.
  \end{greybox}

\begin{tpcfigure}{c03.fig.raeume}
  \demo{c03.demo.raeume}
  \caption{The three spaces side by side, each with the squared moduli of one vector.
    (a) $\mathbb{C}^2$: the vector $\tfrac12\vb{v}$ with $\vb{v}$ from
    Equation~\eqref{c03.eq.c2-vektoren}; its two squared moduli are $\tfrac12$ each.
    (b) $\elltwo$: the squared moduli $\abs{v_j}^2 = 1/j^2$ of the first twelve entries of the
    vector with $v_j = 1/j$; infinitely many more follow. (c) $\Ltwo{\mathbb{R}}$: the squared
    modulus $\abs{f(x)}^2 = e^{-x^2}$ of the function $f(x) = e^{-x^2/2}$. The inner product
    of the vector with itself is a sum of two bars on the left, a sum of infinitely many in the
    middle, and the area under the curve on the right.}
  \end{tpcfigure}
\end{vertiefung}

A complex vector space with an inner product is almost a Hilbert space already. One condition
is still missing, \emph{completeness}. It is mathematically important but subtle, and in this
course it plays no essential role: in $\mathbb{C}^n$ it always holds\zoomref{c03.f3}. Before we
define the Hilbert space we need one more piece of notation, which makes computations with
inner products very much easier: Dirac's notation. The real hurdle at the start of quantum
mechanics is the notation; you have to fight your way through it. At first it is unfamiliar,
in the end it makes everything easier, and it is a one-way street: once you have mastered it,
you do not go back.

% ── Lecture 3 (vl03): Dirac notation, inner products, Hilbert space, bases ────────────────
\section{Dirac notation}\label{c03.s4}

Let $V$ be a complex vector space. In quantum mechanics the following notation is used
universally; that is how it is, and we have to live with it. Mathematicians do not like this
notation, physicists love it: that is the difference between mathematicians and physicists.
From now on we write every vector $\vb{v} \in V$ in the so-called ket notation,
\begin{equation*}
  \vb{v} \mapsto \ket{v} .
\end{equation*}
Inside the ket we drop the bold face: the vertical bar and the angle bracket already say that a
vector is meant.

At this point the notation means nothing more than: ``$\ket{v}$ stands for a vector in the
vector space $V$''. That the two extra symbols bring us an advantage is not obvious. Why we have
introduced this unusual notation will become clear soon.

We work through a few examples so that you become familiar with the notation. There are some
rules of calculation that are only implicit in the notation; you have to have seen them once.

\begin{example}[Kets in $\mathbb{C}^2$]\label{c03.ex.kets}\claim{c03.cl.kets}{derived}
  Let $V = \mathbb{C}^2$, the space with which we shall describe the states of a qubit. Two
  example vectors are
  \begin{equation}\label{c03.eq.e0e1}
    \ket{e_0} = \begin{pmatrix} 1 \\ 0 \end{pmatrix},
    \qquad
    \ket{e_1} = \begin{pmatrix} 0 \\ 1 \end{pmatrix} .
  \end{equation}
  This is how the equation is to be read: $\ket{e_0}$ stands for this column, nothing more. In
  a vector space we may add vectors and multiply them by numbers, and so we form two new
  vectors,
  \begin{equation}\label{c03.eq.epm}
    \ket{e_\pm} = \tfrac{1}{\sqrt{2}}\, \ket{e_0} \pm \tfrac{1}{\sqrt{2}}\, \ket{e_1}
    = \frac{1}{\sqrt{2}} \begin{pmatrix} 1 \\ 0 \end{pmatrix} \pm \frac{1}{\sqrt{2}} \begin{pmatrix} 0 \\ 1 \end{pmatrix}
    = \frac{1}{\sqrt{2}} \begin{pmatrix} 1 \\ \pm 1 \end{pmatrix} .
  \end{equation}
  Which vector is represented, one sees by inserting the columns of
  Equation~\eqref{c03.eq.e0e1}.
\end{example}

% TJO's own forward pointer, v03 [49:20]-[49:46], restored after his answer to G-0072 (a).
The vectors $\ket{e_\pm}$ play an important role in the theory of angular momentum and in the
hydrogen atom; they keep turning up in quantum mechanics, especially in spin systems. In the
coming chapters we shall mostly leave the columns out and prefer to calculate directly
with the notation. At the beginning, however, you should always insert the original columns to
see what is represented; in the exercises you will carry out many such calculations.

By the way, it makes no difference whether we write $\ket{v + w}$ or $\ket{v} + \ket{w}$, and
likewise $\ket{z v}$ stands for $z \ket{v}$ when $z$ is a number:
\begin{equation*}
  \ket{v + w} = \ket{v} + \ket{w} ,
  \qquad
  \ket{z v} = z \ket{v} .
\end{equation*}

The ket notation is complemented by the so-called \emph{bra} vectors. If $\ket{v}$ stands for a
column of $n$ complex numbers, the dagger operation of Section~\ref{c03.s2} turns it into a row,
and for this row we introduce a symbol of its own, $\bra{v}$. When you see this odd symbol, you
have to understand it as a row vector. Bra and ket together write the inner product.

\tpcdef{D81}

\begin{historical}[Bra-ket]
  The notation is due to Paul Dirac, and today it is used universally
  \srcref{tpc-script-v0.2}{vl03, box on the bra-ket notation}. In English a \emph{bracket}: on
  the left stands the ``bra'', on the right the ``ket'', and the ``c'' got lost somewhere.
\end{historical}

This is not a new definition but a new notation for the same objects and the same equations.
For the same number we now have three ways of writing, the sum $\sum_{j} v_j^* w_j$, the product
$\vb{v}\adj\vb{w}$ and the bra-ket $\braket{v}{w}$, and a bra-ket is computed by the old rule.
Between $\braket{v}{w}$ and $(\ket{v})\adj \ket{w}$ we make no distinction at all any more. This
holds for the inner product of Section~\ref{c03.s2}, and we keep using it whenever we insert
columns. In the next section other inner products appear as well. For them the bra $\bra{v}$ is
the assignment that gives every ket $\ket{w}$ the number $\braket{v}{w}$, and in general no
longer the conjugate row (Example~\ref{c03.ex.gewichtet}).

\begin{example}[A bra-ket in $\mathbb{C}^2$]\label{c03.ex.braket}\claim{c03.cl.braket}{derived}
  With the vectors of Example~\ref{c03.ex.kets},
  \begin{equation}\label{c03.eq.e0eplus}
    \braket{e_0}{e_+} = (\ket{e_0})\adj \ket{e_+}
    = \begin{pmatrix} 1 & 0 \end{pmatrix}
      \begin{pmatrix} \frac{1}{\sqrt{2}} \\ \frac{1}{\sqrt{2}} \end{pmatrix}
    = \frac{1}{\sqrt{2}} .
  \end{equation}
\end{example}

For the notation to show its advantage, one needs a rule for how the dagger operation acts on
linear combinations. For complex numbers $\alpha, \beta$ and kets $\ket{v}, \ket{w}$,
\begin{equation}\label{c03.eq.bra-regel}\claim{c03.cl.bra-regel}{derived}
  (\alpha \ket{v} + \beta \ket{w})\adj = \alpha^* \bra{v} + \beta^* \bra{w} ,
\end{equation}
because the dagger operation conjugates every entry, and
$(\alpha v_j + \beta w_j)^* = \alpha^* v_j^* + \beta^* w_j^*$. It is exactly this rule that
makes the notation useful: with it we compute inner products without ever writing the vectors
down as columns. In $\mathbb{C}^2$ writing them down is no great burden. In
infinite-dimensional spaces, however, one generally cannot write a vector down completely at
all, and then calculating with kets and this rule is the way to calculate efficiently.

\begin{example}[An inner product with kets only]\label{c03.ex.dirac}\claim{c03.cl.dirac}{derived}
  We want the inner product $\braket{v}{w}$ of
  \begin{equation}\label{c03.eq.dirac-vw}
    \ket{v} = \begin{pmatrix} 1 \\ i \end{pmatrix},
    \qquad
    \ket{w} = \begin{pmatrix} 1 \\ 1 \end{pmatrix},
  \end{equation}
  this time without inserting the columns. First we express both vectors by the kets of
  Example~\ref{c03.ex.kets}: $\ket{v} = \ket{e_0} + i \ket{e_1}$ and
  $\ket{w} = \ket{e_0} + \ket{e_1}$. By Equation~\eqref{c03.eq.bra-regel},
  $\bra{v} = \bra{e_0} - i \bra{e_1}$; note the minus sign. Multiplying out gives four inner
  products:
  \begin{equation}\label{c03.eq.dirac}
    \begin{aligned}
      \braket{v}{w} &= (\bra{e_0} - i \bra{e_1}) (\ket{e_0} + \ket{e_1}) \\
      &= \braket{e_0}{e_0} + \braket{e_0}{e_1} - i \braket{e_1}{e_0} - i \braket{e_1}{e_1} \\
      &= 1 + 0 - 0 - i = 1 - i ,
    \end{aligned}
  \end{equation}
  because $\braket{e_0}{e_0} = \braket{e_1}{e_1} = 1$ and $\braket{e_0}{e_1} = \braket{e_1}{e_0} = 0$.
  We would have obtained the same result by simply inserting the columns:
  $\braket{v}{w} = \begin{pmatrix} 1 & -i \end{pmatrix} \begin{pmatrix} 1 \\ 1 \end{pmatrix} = 1 - i$.
\end{example}

Let us be honest: in two dimensions the first route is the more laborious one. In higher
dimensions, though, you will very quickly discover a love for the ket notation. I use it in two
dimensions all the same, because it plays together very nicely with operators and matrices; we
shall see that with the operators.

Once you have mastered the ket notation, you can read the equations of quantum mechanics. The
translation is always the same: a ket stands for a vector.

\section{Inner products}\label{c03.s5}

Before we can define the Hilbert space, we must capture the concept of an inner product more
cleanly and mathematically better. In $\mathbb{C}^n$ we have so far always used the inner
product $\vb{v}\adj\vb{w}$. But even on this vector space there are other inner products, in fact
infinitely many (Example~\ref{c03.ex.gewichtet}). In physics one normally chooses one and then works with this
inner product for ever and ever; mathematically there are several possibilities, and sometimes
it helps to be a little flexible. A complex vector space can also exist without any inner
product at all: the inner product is an additional choice, not a structure the vector space
brings along by itself. We need one because quantum mechanics is probabilistic; it is our tool
for probabilities.

An inner product assigns a complex number to any two vectors, and we write it in Dirac's
notation. Not every such assignment counts as an inner product; it has to satisfy a few
conditions. These conditions describe the other inner products as well. For calculations with
columns we keep using the inner product of Section~\ref{c03.s2}.

\tpcdef{D82}

When a number is pulled out of the first argument, a star appears. These rules hold in the
finite-dimensional as in the infinite-dimensional case; that is why the notation and the concept
are so important. The inner product $\vb{v}\adj\vb{w}$ of $\mathbb{C}^n$ satisfies all three
conditions: they are rules~2 to~4 of Proposition~\ref{c03.prop.skalarprodukt}. The complex
conjugation in the second rule for numbers is important so that $\braket{\psi}{\psi}$ can be
positive. If the inner product were linear in the first argument as well, then for every vector
\begin{equation}\label{c03.eq.warum-stern}\claim{c03.cl.warum-stern}{derived}
  \braket{i\psi}{i\psi} = i \cdot i\, \braket{\psi}{\psi} = -\braket{\psi}{\psi} ,
\end{equation}
and both sides could be $\geq 0$ only if $\braket{\psi}{\psi} = 0$, also for
$\ket{\psi} \neq 0$. Example~\ref{c03.ex.ohne-stern} shows the same for a single vector.

\begin{example}[Infinitely many inner products on $\mathbb{C}^2$]\label{c03.ex.gewichtet}\claim{c03.cl.gewichtet}{derived}
  Let $c > 0$ be a real number. For $\vb{v}, \vb{w} \in \mathbb{C}^2$ we set
  \begin{equation}\label{c03.eq.gewichtet}
    \braket{v}{w}_c = v_1^* w_1 + c\, v_2^* w_2 .
  \end{equation}
  For every $c > 0$ this is an inner product: additivity and the two rules for numbers hold
  term by term, exactly as in Proposition~\ref{c03.prop.skalarprodukt}, and
  $\braket{v}{v}_c = \abs{v_1}^2 + c \abs{v_2}^2$ is $\geq 0$ and zero only if
  $v_1 = v_2 = 0$. For $c = 1$ it is the inner product of $\mathbb{C}^2$. Different $c$ give
  different inner products, because $\braket{e_1}{e_1}_c = c$. So there are infinitely many.
  Each comes with its own bra: since
  $\braket{v}{w}_c = \begin{pmatrix} v_1^* & c\, v_2^* \end{pmatrix} \vb{w}$, here $\bra{v}$
  stands for the row $\begin{pmatrix} v_1^* & c\, v_2^* \end{pmatrix}$, and for $c \neq 1$ and
  $v_2 \neq 0$ that is not the row $\vb{v}\adj$.
\end{example}

Positivity allows us to define the length of a vector.

\tpcdef{D83}

In $\mathbb{C}^n$ with the inner product of Section~\ref{c03.s2},
$\norm{v}^2 = \vb{v}\adj\vb{v} = \sum_j \abs{v_j}^2$ (Equation~\eqref{c03.eq.selbst}). Instead of $\norm{\psi}$ one often simply says ``the modulus
of $\psi$''.

Positivity is a very strong condition, and it has two very nice consequences. The first is
conjugate symmetry. For the inner products we have seen so far it is nothing new; here it
follows from the three conditions alone.

\begin{proposition}[Conjugate symmetry]\label{c03.prop.symmetrie}\claim[proof-included]{c03.cl.symmetrie}{derived}
  For every inner product and all vectors $\ket{\phi}, \ket{\psi}$,
  \begin{equation}\label{c03.eq.symmetrie}
    \braket{\phi}{\psi} = \braket{\psi}{\phi}^* .
  \end{equation}
\end{proposition}
\begin{proof}
  By positivity, $\braket{\xi}{\xi}$ is a real number for every vector $\ket{\xi}$. Let
  $a = \braket{\phi}{\psi}$ and $b = \braket{\psi}{\phi}$. By additivity,
  \begin{equation}\label{c03.eq.symmetrie-beweis}
    \braket{\phi + \psi}{\phi + \psi} = \braket{\phi}{\phi} + \braket{\psi}{\psi} + a + b ,
  \end{equation}
  so $a + b$ is real. The same calculation with $i\ket{\psi}$ instead of $\ket{\psi}$: by the
  rules for numbers, $\braket{\phi}{i\psi} = i a$, $\braket{i\psi}{\phi} = -i b$ and
  $\braket{i\psi}{i\psi} = (-i)\, i\, \braket{\psi}{\psi} = \braket{\psi}{\psi}$, so
  $i a - i b = i (a - b)$ is real. From $\operatorname{Im}(a + b) = 0$ it follows that
  $\operatorname{Im} b = -\operatorname{Im} a$, and since $a - b$ is purely imaginary,
  $\operatorname{Re} b = \operatorname{Re} a$. Hence $b = a^*$.
\end{proof}

Reed and Simon include conjugate symmetry as a fourth condition in the definition of the inner
product \srcref{reed-simon-1}{p.~36}; by the proposition this is the same definition.

The second consequence is new, the \emph{Schwarz inequality}. It is particularly beautiful and
particularly useful in quantum mechanics.

\begin{proposition}[Schwarz inequality]\label{c03.prop.schwarz}\claim[proof-included]{c03.cl.schwarz}{derived}
  For every inner product and all vectors $\ket{\phi}, \ket{\psi}$,
  \begin{equation}\label{c03.eq.schwarz}
    \abs{\braket{\phi}{\psi}} \leq \norm{\phi}\, \norm{\psi} .
  \end{equation}
\end{proposition}
\begin{proof}
  For $\ket{\psi} = 0$ both sides are zero. Otherwise $\norm{\psi}^2 > 0$. We set
  $c = \braket{\psi}{\phi} / \norm{\psi}^2$ and $\ket{\xi} = \ket{\phi} - c \ket{\psi}$. By
  Proposition~\ref{c03.prop.symmetrie}, $c^* = \braket{\phi}{\psi} / \norm{\psi}^2$. With
  additivity, the rules for numbers and positivity for $\ket{\xi}$,
  \begin{equation}\label{c03.eq.schwarz-beweis}
    0 \leq \braket{\xi}{\xi}
      = \norm{\phi}^2 - c\, \braket{\phi}{\psi} - c^*\, \braket{\psi}{\phi} + c^* c\, \norm{\psi}^2
      = \norm{\phi}^2 - \frac{\abs{\braket{\phi}{\psi}}^2}{\norm{\psi}^2} ,
  \end{equation}
  because up to its sign each of the last three terms equals the fraction on the right.
  Multiplying by $\norm{\psi}^2$ and taking the square root gives the claim.
\end{proof}

With the Schwarz inequality one can prove another, equally useful inequality, the triangle
inequality: the norm of a sum of vectors is never larger than the sum of the norms.

\begin{proposition}[Triangle inequality]\label{c03.prop.dreieck}\claim[proof-included]{c03.cl.dreieck}{derived}
  For every inner product and finitely many vectors $\ket{\psi_1}, \dots, \ket{\psi_N}$,
  \begin{equation}\label{c03.eq.dreieck}
    \norm{\sum_{j=1}^{N} \ket{\psi_j}} \leq \sum_{j=1}^{N} \norm{\psi_j} .
  \end{equation}
\end{proposition}
\begin{proof}
  For two vectors, by additivity and Proposition~\ref{c03.prop.symmetrie},
  \begin{equation}\label{c03.eq.dreieck-beweis}
    \norm{\phi + \psi}^2
      = \norm{\phi}^2 + \norm{\psi}^2 + 2 \operatorname{Re} \braket{\phi}{\psi}
      \leq \norm{\phi}^2 + \norm{\psi}^2 + 2 \norm{\phi}\, \norm{\psi}
      = \bigl( \norm{\phi} + \norm{\psi} \bigr)^2 ,
  \end{equation}
  because $\braket{\psi}{\phi} = \braket{\phi}{\psi}^*$, $\operatorname{Re} z \leq \abs{z}$ for
  every complex number $z$, and then the Schwarz inequality. The square root gives
  $\norm{\phi + \psi} \leq \norm{\phi} + \norm{\psi}$. For $N$ vectors one applies this
  repeatedly: $\norm{\psi_1 + \dots + \psi_N} \leq \norm{\psi_1 + \dots + \psi_{N-1}} + \norm{\psi_N}$,
  and so on down to a single vector.
\end{proof}

Reed and Simon show in the same way that the norm of a vector space with inner product is a norm
in the sense of normed spaces \srcref{reed-simon-1}{p.~38}.

Both inequalities are working tools that we use practically all the time. The most important
one is the Schwarz inequality: it lies behind Heisenberg's uncertainty relation
\srcref{tpc-script-v0.2}{vl03, section on inner products}, which comes in Chapter~6, through
the idea of its proof. In equation~\eqref{c03.eq.schwarz-beweis} the inner product of a
cleverly chosen vector with itself is never negative; exactly this step also carries the proof
of the uncertainty relation, as Section~\ref{c03.s10} shows. With the triangle inequality we
bound quantities.

One last definition before we leave inner products.

\tpcdef{D84}

This is exactly analogous to ordinary vectors in three-dimensional space. Careful: a vanishing
inner product does \emph{not} mean that one of the vectors is zero. Positivity forces that only
in the case $\ket{\phi} = \ket{\psi}$.

\begin{example}[Orthonormal vectors in $\mathbb{C}^2$]\label{c03.ex.orthogonal}\claim{c03.cl.orthogonal}{derived}
  The vectors $\ket{e_\pm}$ of Example~\ref{c03.ex.kets} are not zero, and yet
  \begin{equation}\label{c03.eq.orthogonal}
    \braket{e_+}{e_-} = \frac{1}{2} \begin{pmatrix} 1 & 1 \end{pmatrix} \begin{pmatrix} 1 \\ -1 \end{pmatrix} = 0 ,
    \qquad
    \norm{e_\pm}^2 = \frac{1}{2} \bigl( 1 + (\pm 1)^2 \bigr) = 1 .
  \end{equation}
  Likewise $\braket{e_0}{e_1} = 0$ and $\norm{e_0} = \norm{e_1} = 1$. Both pairs are
  orthonormal.
\end{example}

One may also read the norm as the length of a vector. The vectors of quantum mechanics live in
an abstract space that has nothing at all to do with our three dimensions; still it is sometimes
helpful to read their norms as lengths and orthogonal vectors as perpendicular to each other.

\section{Hilbert space}\label{c03.s6}

Now we have enough together to define Hilbert spaces.

\tpcdef{D14}

Completeness is a very important but subtle mathematical condition. It has to do with
convergence in infinite-dimensional spaces, and in our lectures it plays no essential role. For
our calculations in $\mathbb{C}^n$ with the inner product of Section~\ref{c03.s2} it holds
automatically\zoomref{c03.f3}. There we can simplify all calculations by ignoring it, and for
these calculations a Hilbert space is simply a complex vector space with an inner product. So
$\mathbb{C}^n$ is a Hilbert space, and for the qubit it is $\mathbb{C}^2$. In infinitely many
dimensions the shortcut does not hold: there completeness is used, for example for the
expansion in a basis with infinitely many vectors.

\section{Bases}\label{c03.s7}

The next concept in the theory of Hilbert spaces is the \emph{basis}. A basis is a set of
vectors with very nice properties: all other vectors of the Hilbert space can be decomposed
along it. We need bases for a practical reason: most calculations in quantum mechanics only
become explicit once a basis has been chosen. Normally one chooses a basis and does all
calculations with respect to this basis.

In the core, the index set $X$ of a basis is finite.

\tpcdef{D85}

Two warnings. Confusingly, the second property is called \emph{completeness}; this is a
different notion of completeness from that of the Hilbert space of Section~\ref{c03.s6}. And $\kron{j}{k}$
is the Kronecker delta, not the Dirac delta function, which comes later.

\begin{example}[Two bases of $\mathbb{C}^2$]\label{c03.ex.basis}\claim{c03.cl.basis}{derived}
  The vectors $\ket{e_0}, \ket{e_1}$ form a basis of $\mathbb{C}^2$, with index set
  $X = \{0, 1\}$. They are orthonormal by Example~\ref{c03.ex.orthogonal}; four inner products
  had to be checked here. Completeness is a little more demanding. Every vector of
  $\mathbb{C}^2$ can be written as $\ket{\psi} = \alpha \ket{e_0} + \beta \ket{e_1}$, with the
  entries $\alpha, \beta \in \mathbb{C}$, and we compute the two sides of the completeness
  condition separately. On the left,
  \begin{equation}\label{c03.eq.basis-links}
    \norm{\psi}^2 = (\alpha^* \bra{e_0} + \beta^* \bra{e_1}) (\alpha \ket{e_0} + \beta \ket{e_1})
    = \abs{\alpha}^2 + \abs{\beta}^2 ,
  \end{equation}
  on the right, with $\braket{e_0}{\psi} = \alpha$ and $\braket{e_1}{\psi} = \beta$,
  \begin{equation}\label{c03.eq.basis-rechts}
    \abs{\braket{e_0}{\psi}}^2 + \abs{\braket{e_1}{\psi}}^2 = \abs{\alpha}^2 + \abs{\beta}^2 .
  \end{equation}
  The two sides agree. In the same way $\ket{e_+}, \ket{e_-}$ form a basis: they are
  orthonormal (Example~\ref{c03.ex.orthogonal}), and with
  $\braket{e_\pm}{\psi} = (\alpha \pm \beta) / \sqrt{2}$,
  \begin{equation}\label{c03.eq.basis-pm}
    \abs{\braket{e_+}{\psi}}^2 + \abs{\braket{e_-}{\psi}}^2
    = \tfrac12 \bigl( \abs{\alpha + \beta}^2 + \abs{\alpha - \beta}^2 \bigr)
    = \abs{\alpha}^2 + \abs{\beta}^2 .
  \end{equation}
\end{example}

With $n$ basis vectors, orthonormality means checking $n^2$ equations; by conjugate symmetry
those with $j \leq k$ suffice, and that is still quite a lot. If you are a little bored, write
down a few vectors and check whether they form a basis. Vectors chosen at random do not in
general form a basis: already $\ket{e_0}$ and $\ket{e_+}$ are not orthogonal, by
Example~\ref{c03.ex.braket}.

The most important property of a basis is that every vector can be expanded in it.

\begin{proposition}[Expansion in a basis]\label{c03.prop.entwicklung}\claim[proof-included]{c03.cl.entwicklung}{derived}
  Let $\{\ket{\phi_j}\}_{j \in X}$ be a basis of $\Hil$ with finite index set $X$. Then for
  every vector $\ket{\Psi} \in \Hil$
  \begin{equation}\label{c03.eq.vollstaendigkeit}
    \ket{\Psi} = \sum_{j \in X} \braket{\phi_j}{\Psi}\, \ket{\phi_j} .
  \end{equation}
  Conversely, an orthonormal system with finite index set in which every vector can be expanded
  in this way is complete, hence a basis.
\end{proposition}
\begin{proof}
  Let $c_j = \braket{\phi_j}{\Psi}$ and $\ket{\xi} = \ket{\Psi} - \sum_j c_j \ket{\phi_j}$. By
  orthonormality, $\braket{\phi_k}{\xi} = c_k - \sum_j c_j \kron{k}{j} = 0$ for every $k$. With
  the rules for numbers in the first argument,
  $\braket{\xi}{\xi} = \braket{\Psi}{\xi} - \sum_j c_j^* \braket{\phi_j}{\xi} = \braket{\Psi}{\xi}$,
  and by Proposition~\ref{c03.prop.symmetrie}, $\braket{\Psi}{\phi_j} = c_j^*$, so
  \begin{equation}\label{c03.eq.entwicklung-beweis}
    \braket{\xi}{\xi} = \braket{\Psi}{\Psi} - \sum_{j} c_j\, \braket{\Psi}{\phi_j}
    = \norm{\Psi}^2 - \sum_{j} \abs{c_j}^2 = 0
  \end{equation}
  by completeness. By positivity $\ket{\xi} = 0$, and that is
  Equation~\eqref{c03.eq.vollstaendigkeit}. Conversely, if the expansion holds, then
  $\norm{\Psi}^2 = \sum_{j,k} c_j^* c_k \braket{\phi_j}{\phi_k} = \sum_j \abs{c_j}^2$.
\end{proof}

\begin{merksatz}[Completeness relation]
  The equation
  $\ket{\Psi} = \sum_{j \in X} \braket{\phi_j}{\Psi}\, \ket{\phi_j}$ is one form of the
  so-called \emph{completeness relation}: every vector is the sum of its components along the
  basis vectors, and the coefficients are inner products. It looks harmless, but it is used all
  the time in quantum mechanics.
\end{merksatz}

\begin{example}[One vector, two bases]\label{c03.ex.zwei-basen}\claim{c03.cl.zwei-basen}{derived}
  Let $\ket{\Psi} = \tfrac{1}{\sqrt{5}} \bigl( 2 \ket{e_0} + \ket{e_1} \bigr)$. In the basis
  $\ket{e_0}, \ket{e_1}$ the coefficients are $c_0 = \braket{e_0}{\Psi} = 2/\sqrt{5}$ and
  $c_1 = \braket{e_1}{\Psi} = 1/\sqrt{5}$, in the basis $\ket{e_+}, \ket{e_-}$ they are
  $c_\pm = \braket{e_\pm}{\Psi} = (2 \pm 1)/\sqrt{10}$. The weights
  \begin{equation}\label{c03.eq.gewichte}
    \abs{c_0}^2 = \tfrac45 , \quad \abs{c_1}^2 = \tfrac15 ,
    \qquad
    \abs{c_+}^2 = \tfrac{9}{10} , \quad \abs{c_-}^2 = \tfrac{1}{10}
  \end{equation}
  differ between the two bases; their sum is $1 = \norm{\Psi}^2$ in both.
\end{example}

Figure~\ref{c03.fig.vollst} works this out for the two bases of $\mathbb{C}^2$. In (a) each of
the two decompositions of the same $\ket{\Psi}$ closes up into a rectangle of its own: different
sides, the same diagonal. Panel (b) shows what stays independent of the basis: not the
individual weights, but their sum.

\begin{tpcfigure}{c03.fig.vollst}
  \demo{c03.demo.vollst}
  \caption{The completeness relation for one vector in two bases. (a) The vector
    $\ket{\Psi} = (2\ket{e_0} + \ket{e_1})/\sqrt{5}$ of Example~\ref{c03.ex.zwei-basen} (thick
    arrow) as the sum of its components $c_j \ket{\phi_j}$ with $c_j = \braket{\phi_j}{\Psi}$,
    once in the basis $\ket{e_0}, \ket{e_1}$ (rectangle closed with dashes), once in the basis
    $\ket{e_+}, \ket{e_-}$ (rectangle closed with dots). Since all coefficients are real,
    everything lies in the real plane spanned by $\ket{e_0}$ and $\ket{e_1}$; the rays of
    $\ket{e_\pm}$ bisect the angles between the axes. The two rectangles have different sides
    and the same diagonal. (b) The weights $\abs{c_j}^2$ as bar segments: $\tfrac45$ and
    $\tfrac15$ against $\tfrac{9}{10}$ and $\tfrac{1}{10}$. The split depends on the basis, the
    sum $\norm{\Psi}^2 = 1$ (dashed line) does not.}
\end{tpcfigure}

In finitely many dimensions it is relatively easy to find a basis. In $\mathbb{C}^n$ with the
inner product of Section~\ref{c03.s2} one simply takes the columns $\ket{u_1}, \dots, \ket{u_n}$, where $\ket{u_j}$ has a $1$ in the $j$-th entry
and zeros everywhere else; for $n = 2$ these are $\ket{e_0}$ and $\ket{e_1}$. For a vector
$\ket{\psi}$ with entries $\psi_1, \dots, \psi_n$,
\begin{equation}\label{c03.eq.kanonisch}\claim{c03.cl.kanonisch}{derived}
  \braket{u_j}{u_k} = \kron{j}{k} ,
  \qquad
  \sum_{j=1}^{n} \abs{\braket{u_j}{\psi}}^2 = \sum_{j=1}^{n} \abs{\psi_j}^2 = \norm{\psi}^2 .
\end{equation}
In the infinite-dimensional case everything is harder: that a basis exists is not clear, it has
to be proved, and the proof is hard\zoomref{c03.v4}.

\begin{vertiefung}[id=c03.v4,
    titel={Bases with infinitely many vectors},
    wozu={The bases of $\elltwo$ and $\Ltwo{\mathbb{R}}$ in this block have infinitely many
      vectors, and the sums of the completeness relation become series. Here is how a basis is
      defined then, that there always is one and that the expansion in a basis with a countable
      index set holds, with a basis of $\elltwo$ and one of $\Ltwo{\mathbb{R}}$.},
    voraussetzung={Section~\ref{c03.s7} up to Proposition~\ref{c03.prop.entwicklung}; blocks~\ref{c03.v1}
      and~\ref{c03.v2}; the completeness of $\elltwo$ and
      $\Ltwo{\mathbb{R}}$ from Foundations~\ref{c03.f3}. The core up to
      Chapter~4 does not need this block.},
    zurueck={In infinite dimension, too, every Hilbert space has a basis. If its index set is
      countable, every vector is the sum of its components, as a series that converges in
      norm. In the core the finite sum suffices.}]
  From now on we assume that we always have a basis, otherwise it gets too hard. For Hilbert
  spaces this is more than an assumption, namely a theorem. First, however, one has to say
  what a basis with infinitely many vectors is. For a countable index set, such as
  $X = \{1, 2, 3, \dots\}$, the sums become series, and the completeness of \defref{D85} then
  requires series to converge. But an index set may also be uncountable. Reed and Simon
  therefore define the basis without a sum:

  \tpcdef{D86}

  \begin{theorem}[Existence of a basis]\label{c03.thm.existenz}\claim[proof-omitted]{c03.cl.existenz}{source-quoted}
    Every Hilbert space has an orthonormal basis.
  \end{theorem}

  Reed and Simon prove this with Zorn's lemma \srcref{reed-simon-1}{pp.~44--45}. The next
  theorem says that the expansion in a basis also holds when it has infinitely many vectors,
  countably many.

  \begin{theorem}[Expansion in an orthonormal basis]\label{c03.thm.parseval}\claim[proof-omitted]{c03.cl.parseval}{source-quoted}
    Let $\{\ket{\phi_j}\}_{j \in X}$ be an orthonormal basis of a Hilbert space $\Hil$ with a
    countable index set $X$. Then for every vector $\ket{\Psi} \in \Hil$
    \begin{equation}\label{c03.eq.parseval}
      \ket{\Psi} = \sum_{j \in X} \braket{\phi_j}{\Psi}\, \ket{\phi_j} ,
      \qquad
      \norm{\Psi}^2 = \sum_{j \in X} \abs{\braket{\phi_j}{\Psi}}^2 .
    \end{equation}
    The first sum is a series of vectors; it converges in norm, independent of the order of the
    terms, to $\ket{\Psi}$. Conversely, for complex numbers $c_j$ with
    $\sum_{j \in X} \abs{c_j}^2 < \infty$ the sum $\sum_{j \in X} c_j \ket{\phi_j}$ converges in
    norm to a vector in $\Hil$.
  \end{theorem}

  This is Theorem~II.6 of Reed and Simon, in the case of a countable index set; there the
  second equation is called Parseval's relation \srcref{reed-simon-1}{pp.~45--46}. Reed and
  Simon state the theorem for every index set. We need only countable ones: the bases of
  $\elltwo$ and $\Ltwo{\mathbb{R}}$ below are enumerated by $k = 1, 2, 3, \dots$ and
  $n = 0, 1, 2, \dots$. So the definition of the core fits the one of Reed and Simon:

  \begin{proposition}[Two definitions, one basis]\label{c03.prop.maximal}\claim[proof-included]{c03.cl.maximal}{derived}
    An orthonormal system $\{\ket{\phi_j}\}_{j \in X}$ of a Hilbert space $\Hil$ with a
    countable index set $X$ is an orthonormal basis in the sense of \defref{D86} if and only if
    $\norm{\psi}^2 = \sum_{j \in X} \abs{\braket{\phi_j}{\psi}}^2$ for every vector
    $\ket{\psi} \in \Hil$. For finite $X$ this is the basis of \defref{D85}.
  \end{proposition}
  \begin{proof}
    If the system is an orthonormal basis, the equation holds by
    Theorem~\ref{c03.thm.parseval}. If it is not, there is an orthonormal system that contains
    all the $\ket{\phi_j}$ and one further vector $\ket{\xi}$. Then $\norm{\xi} = 1$ and
    $\braket{\phi_j}{\xi} = 0$ for every $j$; every finite partial sum of
    $\sum_j \abs{\braket{\phi_j}{\xi}}^2$ is zero, and the equation for $\ket{\psi} = \ket{\xi}$
    would say $1 = 0$.
  \end{proof}

  \begin{example}[A basis of $\elltwo$]\label{c03.ex.elltwo-basis}\claim{c03.cl.elltwo-basis}{derived}
    Let $\vb{u}_k \in \elltwo$ be the vector with a $1$ in the $k$-th entry and zeros everywhere
    else, $k = 1, 2, 3, \dots$ Then $\vb{u}_j\adj\vb{u}_k = \kron{j}{k}$, and for every
    $\vb{v} \in \elltwo$, $\vb{u}_k\adj\vb{v} = v_k$, so
    \begin{equation}\label{c03.eq.elltwo-basis}
      \sum_{k=1}^{\infty} \abs{\vb{u}_k\adj\vb{v}}^2 = \sum_{k=1}^{\infty} \abs{v_k}^2 = \vb{v}\adj\vb{v} .
    \end{equation}
    Since $\elltwo$ is a Hilbert space by Theorem~\ref{c03.thm.vollstaendig}, the $\vb{u}_k$
    form an orthonormal basis of $\elltwo$ by Proposition~\ref{c03.prop.maximal}. The
    completeness relation says here: the partial sums $\sum_{k \leq N} v_k \vb{u}_k$, that is
    the first $N$ entries of $\vb{v}$ followed by zeros, converge to $\vb{v}$. One also sees this
    directly: the squared distance is $\sum_{k > N} \abs{v_k}^2$, the remainder of a convergent
    series.
  \end{example}

  For $\Ltwo{\mathbb{R}}$, Theorem~\ref{c03.thm.existenz} gives a basis, because the space in
  the version of Foundations~\ref{c03.f2} is a Hilbert space by
  Theorem~\ref{c03.thm.vollstaendig}. Reed and Simon also name a concrete one.

  \begin{theorem}[Hermite functions]\label{c03.thm.hermite}\claim[proof-omitted]{c03.cl.hermite}{source-quoted}
    The Hermite functions $\phi_n$, $n = 0, 1, 2, \dots$, form an orthonormal basis of
    $\Ltwo{\mathbb{R}}$. The first one is $\phi_0(x) = \pi^{-1/4}\, e^{-x^2/2}$.
  \end{theorem}

  Reed and Simon also call them the wave functions of the harmonic oscillator and refer to
  exercises for the proof \srcref{reed-simon-1}{p.~142}; they come back in Chapter~7. That
  $\phi_0$ is normalised follows from Equation~\eqref{c03.eq.gauss}:
  \begin{equation}\label{c03.eq.phinull}\claim{c03.cl.phinull}{derived}
    \braket{\phi_0}{\phi_0} = \pi^{-1/2} \int_{-\infty}^{\infty} e^{-x^2}\, \dd x
    = \pi^{-1/2} \sqrt{\pi} = 1 .
  \end{equation}
\end{vertiefung}

\subsection{The Gram-Schmidt procedure}\label{c03.s7.1}

How does one build bases in practice? With the so-called orthonormalisation procedure of Gram
and Schmidt. One starts with some list of vectors $\ket{\Phi_1}, \ket{\Phi_2}, \dots$; these
vectors are completely arbitrary and need not satisfy any property. From them one determines,
one after the other, pairwise orthogonal vectors $\ket{\Psi_n}$, with a single formula.

\begin{proposition}[Gram-Schmidt procedure]\label{c03.prop.gram-schmidt}\claim[proof-included]{c03.cl.gram-schmidt}{derived}
  Let $\ket{\Phi_1}, \dots, \ket{\Phi_N}$ be vectors of a Hilbert space $\Hil$, and let
  $\ket{\Psi_1}, \dots, \ket{\Psi_N}$ be determined one after the other by
  \begin{equation}\label{c03.eq.gram-schmidt}
    \ket{\Psi_n} = \ket{\Phi_n} - \sum_{k=1}^{n-1} \ket{\Psi_k}\, \frac{\braket{\Psi_k}{\Phi_n}}{\norm{\Psi_k}^2}
  \end{equation}
  where terms with $\ket{\Psi_k} = 0$ are left out. Then:
  \begin{enumerate}
    \item $\braket{\Psi_k}{\Psi_n} = 0$ for $k < n$.
    \item For every $m \leq N$, every $\ket{\Phi_n}$ with $n \leq m$ is a linear combination of
      the $\ket{\Psi_k}$ with $k \leq m$, and conversely.
    \item The vectors $\ket{\Psi_n} / \norm{\Psi_n}$ with $\ket{\Psi_n} \neq 0$ form an
      orthonormal system.
    \item If only the zero vector is orthogonal to all the $\ket{\Phi_n}$, this system is a
      basis of $\Hil$.
  \end{enumerate}
\end{proposition}
\begin{proof}
  Parts~1 and~2 by induction over $n$; for $n = 1$, $\ket{\Psi_1} = \ket{\Phi_1}$. Let
  $\ket{\Psi_1}, \dots, \ket{\Psi_{n-1}}$ be pairwise orthogonal. For $j < n$ with
  $\ket{\Psi_j} \neq 0$, in
  \begin{equation}\label{c03.eq.gs-beweis}
    \braket{\Psi_j}{\Psi_n}
    = \braket{\Psi_j}{\Phi_n} - \sum_{k} \braket{\Psi_j}{\Psi_k}\, \frac{\braket{\Psi_k}{\Phi_n}}{\norm{\Psi_k}^2}
    = \braket{\Psi_j}{\Phi_n} - \braket{\Psi_j}{\Phi_n} = 0
  \end{equation}
  only the term $k = j$ of the sum survives; for $\ket{\Psi_j} = 0$ the inner product is zero
  anyway. By Equation~\eqref{c03.eq.gram-schmidt}, $\ket{\Phi_n}$ equals $\ket{\Psi_n}$ plus a
  linear combination of the $\ket{\Psi_k}$ with $k < n$, and $\ket{\Psi_n}$ is $\ket{\Phi_n}$
  minus a linear combination of the $\ket{\Psi_k}$ with $k < n$, which by the induction
  hypothesis are linear combinations of the $\ket{\Phi_k}$ with $k < n$. Part~3: by the rules for
  numbers, $\norm{z\psi} = \abs{z}\, \norm{\psi}$, so the divided vectors have norm $1$, and
  they are orthogonal by part~1. Part~4: let $\ket{\xi_1}, \dots, \ket{\xi_M}$ be the vectors of
  part~3, let $\ket{\psi} \in \Hil$ and
  $\ket{r} = \ket{\psi} - \sum_m \braket{\xi_m}{\psi}\, \ket{\xi_m}$. As in the proof of
  Proposition~\ref{c03.prop.entwicklung}, $\braket{\xi_m}{r} = 0$ for every $m$, hence also
  $\braket{\Psi_k}{r} = 0$ for every $k$ and, by part~2, $\braket{\Phi_n}{r} = 0$ for every
  $n$. By the hypothesis, $\ket{r} = 0$. So every vector can be expanded in the $\ket{\xi_m}$,
  and by Proposition~\ref{c03.prop.entwicklung} the system is complete.
\end{proof}

All quantities in the formula can be computed: the inner product is a complex number, the norm
a nonnegative real number, the $\ket{\Psi_k}$ are vectors, and the sum is a sum of vectors. Here the number stands to
the right of the vector; in the definition it always stood on the left. I do that very often,
and you have to live with it: what is meant is $\ket{\Psi_k}\, z = z \ket{\Psi_k}$. The index $n$
has nothing at all to do with the dimension of the space; it only counts off the list. The terms
with $\ket{\Psi_k} = 0$ have to be left out, because otherwise $\norm{\Psi_k}^2 = 0$ would stand
in the denominator and we would be dividing by zero. It may well happen that some $\ket{\Psi_n}$
are zero (Exercise~\ref{c03.exc.gram-schmidt}).

Because the procedure consists of several parts and can be confusing, we work through it once
completely.

\begin{example}[Gram-Schmidt in $\mathbb{C}^2$]\label{c03.ex.gram-schmidt}\claim{c03.cl.gs-beispiel}{derived}
  We work in my favourite Hilbert space $\mathbb{C}^2$ and start with two vectors that I have
  just made up:
  \begin{equation}\label{c03.eq.gs-start}
    \ket{\Phi_1} = \begin{pmatrix} 2 \\ 0 \end{pmatrix},
    \qquad
    \ket{\Phi_2} = \begin{pmatrix} 1 \\ 1 \end{pmatrix} .
  \end{equation}
  In the first step the sum runs from $1$ to $0$; so nothing is summed at all, and
  $\ket{\Psi_1} = \ket{\Phi_1}$. In the second step only the term $k = 1$ remains. With
  $\braket{\Psi_1}{\Phi_2} = 2$ and $\norm{\Psi_1}^2 = 4$,
  \begin{equation}\label{c03.eq.gs-beispiel}
    \ket{\Psi_2} = \ket{\Phi_2} - \ket{\Psi_1}\, \frac{\braket{\Psi_1}{\Phi_2}}{\norm{\Psi_1}^2}
    = \begin{pmatrix} 1 \\ 1 \end{pmatrix} - \begin{pmatrix} 2 \\ 0 \end{pmatrix} \frac{2}{4}
    = \begin{pmatrix} 0 \\ 1 \end{pmatrix} .
  \end{equation}
  Divided by the norms $\norm{\Psi_1} = 2$ and $\norm{\Psi_2} = 1$, the made-up list turns into
  the basis $\ket{e_0}, \ket{e_1}$. The condition of part~4 of the proposition holds: from
  $\braket{\Phi_1}{\psi} = 2\psi_1 = 0$ and $\braket{\Phi_2}{\psi} = \psi_1 + \psi_2 = 0$ it
  follows that $\ket{\psi} = 0$.
\end{example}

In higher dimensions the calculation goes the same way, it only gets more and more laborious.
The procedure always works, in the sense that it always yields an orthonormal system; that
system is a basis when the condition of part~4 holds.

\begin{exercise}[A vector drops out]\label{c03.exc.gram-schmidt}\claim{c03.cl.gs-null}{derived}
  Add the third vector $\ket{\Phi_3} = \begin{pmatrix} 1 \\ -1 \end{pmatrix}$ to the two vectors
  of Example~\ref{c03.ex.gram-schmidt} and compute $\ket{\Psi_3}$.
\end{exercise}
\begin{solution}
  With $\ket{\Psi_1} = 2\ket{e_0}$, $\ket{\Psi_2} = \ket{e_1}$, $\braket{\Psi_1}{\Phi_3} = 2$,
  $\norm{\Psi_1}^2 = 4$, $\braket{\Psi_2}{\Phi_3} = -1$ and $\norm{\Psi_2}^2 = 1$,
  \begin{equation}\label{c03.eq.gs-null}
    \ket{\Psi_3} = \begin{pmatrix} 1 \\ -1 \end{pmatrix} - \begin{pmatrix} 2 \\ 0 \end{pmatrix} \frac{2}{4}
    - \begin{pmatrix} 0 \\ 1 \end{pmatrix} \frac{-1}{1} = \begin{pmatrix} 0 \\ 0 \end{pmatrix} .
  \end{equation}
  The third vector is already a linear combination of the first two,
  $\ket{\Phi_3} = \ket{\Phi_1} - \ket{\Phi_2}$, and drops out; the orthonormal system remains
  $\ket{e_0}, \ket{e_1}$.
\end{solution}

% ── Lecture 4 (vl04, with the vl03 subsection on operators, vl03.tex:253-277): operators ──
% Voice loci and the reasons for what is left out are commented in de.tex.
\section{Operators}\label{c03.s8}

We have discussed the important properties of the Hilbert space; now we start with operators,
the second pillar of the theory. Quantum mechanics does not get by with vectors alone: its
states, its measurements and its dynamics are described by operators, as
Section~\ref{c03.s10} shows in overview. The word ``operator'' sounds complicated, but in the
finite-dimensional case an operator is nothing but a matrix. Because we shall also have
infinite-dimensional Hilbert spaces, we nevertheless use the more general word: there, there
is no finite matrix any more; one would have infinitely many rows and columns.

\tpcdef{D87}

As soon as a map has these two properties, we call it an operator. In the core, $\Hil$ is
finite-dimensional, that is: $\Hil$ has a basis with a finite index set (\defref{D85}), for
example $\mathbb{C}^n$ with the inner product of Section~\ref{c03.s2}.

\begin{example}[Operators on $\mathbb{C}^n$ are matrices]\label{c03.ex.matrix}\claim{c03.cl.matrix}{derived}
  Let $\Hil = \mathbb{C}^2$, so the dimension is $n = 2$. Every $2 \times 2$ matrix with
  complex entries $a_{jk}$ is an operator on $\mathbb{C}^2$:
  \begin{equation}\label{c03.eq.matrix}
    \op{A} = \begin{pmatrix} a_{11} & a_{12} \\ a_{21} & a_{22} \end{pmatrix} ,
    \qquad
    \op{A}\ket{\psi} = \begin{pmatrix} a_{11}\psi_1 + a_{12}\psi_2 \\ a_{21}\psi_1 + a_{22}\psi_2 \end{pmatrix} .
  \end{equation}
  The rules for multiplying matrices and vectors correspond exactly to the two conditions of
  the definition, because every entry of $\op{A}\ket{\psi}$ is a sum of multiples of the
  entries $\psi_k$. Conversely, every operator on $\mathbb{C}^2$ is such a matrix: with the
  canonical columns $\ket{u_1}, \ket{u_2}$ of Section~\ref{c03.s7} we have
  $\ket{\psi} = \psi_1 \ket{u_1} + \psi_2 \ket{u_2}$, hence
  $\op{A}\ket{\psi} = \psi_1\, \op{A}\ket{u_1} + \psi_2\, \op{A}\ket{u_2}$, and the columns of
  the matrix are the vectors $\op{A}\ket{u_1}$ and $\op{A}\ket{u_2}$. For every $n$ it works
  the same way, with $n \times n$ matrices.
\end{example}

So it does no harm to think of all operators as matrices: when you hear the word ``operator'',
think of a table. In the infinite-dimensional case this table runs on without end in both
directions, like an Excel sheet that never stops. Then one must be careful, though: the table
belongs to a chosen basis, its sums become series whose convergence must be checked, and some
operators are defined only on part of the space (Section~\ref{c03.s11} and Chapter~5).

\subsection{Identity operator, null operator and the hat}\label{c03.s8.1}

There are a few operators that are trivial but very useful as building blocks.

\tpcdef{D88}

With the null operator we have introduced the notation that we use from now on throughout the
course. With it we avoid most of the confusion.

\begin{merksatz}[Notation]
  A \emph{hat} always means an operator, $\op{A}$; the \emph{bar} of the ket notation a vector,
  $\ket{\psi}$; a symbol without decoration a number. The exceptions are the identity operator
  $\Id$ and two kinds of operators that appear only later: the states (Chapter~4) and the Pauli
  matrices. They carry no hat.
\end{merksatz}

\section{The adjoint operator}\label{c03.s9}

The dagger operation of Section~\ref{c03.s2} turns a column into a row. We can also apply it to
operators, by defining the so-called adjoint operator.

\tpcdef{D17}

That this equation really gives a definition is not self-evident: there must be exactly one
operator that satisfies it.

\begin{proposition}[Existence and uniqueness]\label{c03.prop.adjungiert}\claim[proof-included]{c03.cl.adjungiert}{derived}
  Let $\{\ket{\phi_x}\}_{x \in X}$ be a basis of $\Hil$ with a finite index set $X$, and let
  $\op{A}$ be an operator on $\Hil$. Then there is exactly one operator $\op{A}\adj$ with
  $\braket{\op{A}\adj\phi}{\psi} = \braket{\phi}{\op{A}\psi}$ for all
  $\ket{\phi}, \ket{\psi} \in \Hil$, namely
  \begin{equation}\label{c03.eq.adjungiert}
    \op{A}\adj\ket{\phi} = \sum_{x \in X} \braket{\op{A}\phi_x}{\phi}\, \ket{\phi_x} .
  \end{equation}
\end{proposition}
\begin{proof}
  Uniqueness: if $\op{B}$ and $\op{B}'$ both satisfy the equation, then by additivity
  $\braket{\op{B}\phi - \op{B}'\phi}{\psi} = 0$ for all $\ket{\psi}$. For
  $\ket{\psi} = \op{B}\ket{\phi} - \op{B}'\ket{\phi}$ this is the squared norm of that vector;
  by positivity, therefore, $\op{B}\ket{\phi} = \op{B}'\ket{\phi}$. Existence: the right-hand
  side of Equation~\eqref{c03.eq.adjungiert} depends linearly on $\ket{\phi}$, because the inner
  product is linear in the second argument; so it defines an operator $\op{B}$. By the rules for
  numbers in the first argument and Proposition~\ref{c03.prop.symmetrie},
  \begin{equation}\label{c03.eq.adjungiert-beweis}
    \braket{\op{B}\phi}{\psi}
    = \sum_{x \in X} \braket{\op{A}\phi_x}{\phi}^*\, \braket{\phi_x}{\psi}
    = \sum_{x \in X} \braket{\phi}{\op{A}\phi_x}\, \braket{\phi_x}{\psi}
    = \braket{\phi}{\op{A}\psi} .
  \end{equation}
  In the last step the numbers $\braket{\phi_x}{\psi}$ have moved, by the linearity of $\op{A}$
  and of the inner product, into the second argument, where
  $\sum_x \braket{\phi_x}{\psi}\, \ket{\phi_x} = \ket{\psi}$ stands
  (Proposition~\ref{c03.prop.entwicklung}).
\end{proof}

In the finite-dimensional case the operation can be described explicitly. In $\mathbb{C}^n$
with the inner product of Section~\ref{c03.s2} the dagger means: first transpose, then complex
conjugate all entries,
\begin{equation}\label{c03.eq.adj-matrix}\claim{c03.cl.adj-matrix}{derived}
  \begin{pmatrix}
    a_{11} & a_{12} & \cdots & a_{1n} \\
    a_{21} & a_{22} & \cdots & a_{2n} \\
    \vdots & \vdots & \ddots & \vdots \\
    a_{n1} & a_{n2} & \cdots & a_{nn}
  \end{pmatrix}\adj
  =
  \begin{pmatrix}
    a_{11}^* & a_{21}^* & \cdots & a_{n1}^* \\
    a_{12}^* & a_{22}^* & \cdots & a_{n2}^* \\
    \vdots & \vdots & \ddots & \vdots \\
    a_{1n}^* & a_{2n}^* & \cdots & a_{nn}^*
  \end{pmatrix} .
\end{equation}
For
$\braket{\phi}{\op{A}\psi} = \sum_{j,k} \phi_j^*\, a_{jk}\, \psi_k = \sum_k \bigl( \sum_j a_{jk}^* \phi_j \bigr)^* \psi_k$,
and $\sum_j a_{jk}^* \phi_j$ is the $k$-th entry of the matrix on the right applied to
$\ket{\phi}$. The diagonal entries stay where they are under transposition and are only
conjugated. Like the bra, the adjoint operator depends on the inner product. For the inner
product $\braket{v}{w}_c$ of Example~\ref{c03.ex.gewichtet} and
$\op{A} = \begin{pmatrix} 0 & 1 \\ 0 & 0 \end{pmatrix}$ it is
$\begin{pmatrix} 0 & 0 \\ 1/c & 0 \end{pmatrix}$, as one checks by inserting into the
definition, and for $c \neq 1$ that is not the conjugate transposed matrix.

The product of two operators occurs in the rules of computation. The product $\op{A}\op{B}$ is
the composition: first $\op{B}$, then $\op{A}$, that is
$(\op{A}\op{B})\ket{\psi} = \op{A}(\op{B}\ket{\psi})$; for matrices it is the matrix product.
In general $\op{A}\op{B} \neq \op{B}\op{A}$, for example
\begin{equation}\label{c03.eq.produkt}\claim{c03.cl.produkt}{derived}
  \begin{pmatrix} 0 & 1 \\ 0 & 0 \end{pmatrix} \begin{pmatrix} 0 & 0 \\ 1 & 0 \end{pmatrix}
  = \begin{pmatrix} 1 & 0 \\ 0 & 0 \end{pmatrix} ,
  \qquad
  \begin{pmatrix} 0 & 0 \\ 1 & 0 \end{pmatrix} \begin{pmatrix} 0 & 1 \\ 0 & 0 \end{pmatrix}
  = \begin{pmatrix} 0 & 0 \\ 0 & 1 \end{pmatrix} .
\end{equation}

We shall use the rules of computation in \defref{D17} all the time. You can prove them yourself,
and all of them in the same way.

\begin{proposition}[Rules of computation]\label{c03.prop.adj-regeln}\claim[proof-included]{c03.cl.adj-regeln}{derived}
  On a finite-dimensional Hilbert space the four rules of \defref{D17} hold for all operators
  $\op{A}, \op{B}$ and all $z \in \mathbb{C}$.
\end{proposition}
\begin{proof}
  By Proposition~\ref{c03.prop.adjungiert} it suffices to check the defining equation for each
  right-hand side. For all $\ket{\phi}, \ket{\psi}$, by additivity and the rules for numbers,
  \begin{equation}\label{c03.eq.adj-regeln}
    \begin{aligned}
      \braket{(\op{A}\adj + \op{B}\adj)\phi}{\psi}
        &= \braket{\phi}{\op{A}\psi} + \braket{\phi}{\op{B}\psi}
         = \braket{\phi}{(\op{A} + \op{B})\psi} , \\
      \braket{z^* \op{A}\adj\phi}{\psi}
        &= z\, \braket{\op{A}\adj\phi}{\psi}
         = \braket{\phi}{z\op{A}\psi} , \\
      \braket{\op{B}\adj\op{A}\adj\phi}{\psi}
        &= \braket{\op{A}\adj\phi}{\op{B}\psi}
         = \braket{\phi}{\op{A}\op{B}\psi} .
    \end{aligned}
  \end{equation}
  For the fourth rule we show that $\op{A}$ satisfies the equation for the adjoint of
  $\op{A}\adj$. By Proposition~\ref{c03.prop.symmetrie},
  $\braket{\op{A}\phi}{\psi} = \braket{\psi}{\op{A}\phi}^* = \braket{\op{A}\adj\psi}{\phi}^* = \braket{\phi}{\op{A}\adj\psi}$.
\end{proof}

\section{Classes of operators}\label{c03.s10}

With the identity operator and the adjoint operator alone we can define enormously important
classes of operators. There are six.

\tpcdef{D18}

\tpcdef{D20}

\tpcdef{D89}

\tpcdef{D90}

\tpcdef{D91}

\tpcdef{D92}

All six classes play a very specific role in quantum mechanics, which we shall get to know in
the following chapters \srcref{tpc-script-v0.2}{vl04, paragraph after the list of classes}:
\begin{equation}\label{c03.eq.rollen}\claim{c03.cl.rollen}{source-backed}
  \begin{aligned}
    &\text{\tpcifenglish{Hermitian}{hermitesch}} && \text{\tpcifenglish{observables}{Observablen}} \\
    &\text{\tpcifenglish{positive}{positiv}} && \text{\tpcifenglish{states (with a normalisation, Chapter 4)}{Zustände (mit einer Normierung, Kapitel 4)}} \\
    &\text{\tpcifenglish{projector}{Projektor}} && \text{\tpcifenglish{yes-no measurements (event or no event)}{Ja-Nein-Messungen (Ereignis oder kein Ereignis)}} \\
    &\text{\tpcifenglish{isometric, unitary}{isometrisch, unitär}} && \text{\tpcifenglish{dynamics}{Dynamik}}
  \end{aligned}
\end{equation}
The normal operators are a
wider class in which the others lie (Proposition~\ref{c03.prop.isometrie}).

Hermitian operators on a finite-dimensional space are also called \emph{self-adjoint}; there the
two words mean the same. For operators that are defined only on part of the space one has to
look more closely\zoomref{c03.v5}.

\begin{vertiefung}[id=c03.v5,
    titel={Bounded and unbounded operators},
    wozu={In the core every operator is defined on the whole space, and Hermitian means the same
      as self-adjoint. Here is why this is no longer so for operators with a restricted domain:
      important operators are unbounded in infinitely many dimensions, and an unbounded
      symmetric operator cannot be defined on the whole space.},
    voraussetzung={Sections~\ref{c03.s8} to~\ref{c03.s10}; Deeper blocks~\ref{c03.v1} and~\ref{c03.v4}
      (the example of a basis of $\elltwo$). The core up to Chapter~4 does not need this block;
      Chapter~5 takes the subject up again with domains.},
    zurueck={In finite dimension every operator is bounded and defined on the whole space, and
      Hermitian and self-adjoint are the same. In infinitely many dimensions an operator comes
      with its domain, and only with it can one say whether a symmetric operator is
      self-adjoint.}]
  An operator $\op{T}$ is called \emph{bounded} if there is a number $C \geq 0$ with
  $\norm{\op{T}\psi} \leq C\, \norm{\psi}$ for all vectors $\ket{\psi}$; Reed and Simon define it
  this way for linear maps between normed spaces \srcref{reed-simon-1}{p.~8}. On
  $\mathbb{C}^n$ every operator is bounded. If $\op{A}$ is the matrix with entries $a_{jk}$,
  then by the Schwarz inequality in $\mathbb{C}^n$ we have for every entry
  $\abs{\sum_k a_{jk} \psi_k}^2 \leq \sum_k \abs{a_{jk}}^2\, \norm{\psi}^2$, and the sum over $j$
  gives
  \begin{equation}\label{c03.eq.beschraenkt}\claim{c03.cl.beschraenkt}{derived}
    \norm{\op{A}\psi}^2 \leq \Bigl( \sum_{j,k=1}^{n} \abs{a_{jk}}^2 \Bigr) \norm{\psi}^2 .
  \end{equation}
  In $\elltwo$ this is different.

  \begin{example}[An unbounded operator]\label{c03.ex.unbeschraenkt}\claim{c03.cl.unbeschraenkt}{derived}
    Let $\op{M}$ be the map that multiplies the $j$-th entry of a vector by $j$,
    $(\op{M}\vb{v})_j = j\, v_j$. For the basis vectors $\vb{u}_k$ of
    Example~\ref{c03.ex.elltwo-basis}, $\op{M}\vb{u}_k = k\, \vb{u}_k$, so
    $\norm{\op{M}\vb{u}_k} = k\, \norm{\vb{u}_k}$, and no number $C$ is larger than every $k$.
    Worse: the vector with $v_j = 1/j$ lies in $\elltwo$ (Example~\ref{c03.ex.rate}), but
    $(\op{M}\vb{v})_j = 1$ for every $j$, and the vector of all ones does not lie in $\elltwo$
    (Example~\ref{c03.ex.einsen}). So $\op{M}$ cannot be applied to every vector of $\elltwo$,
    only to those with $\sum_j j^2 \abs{v_j}^2 < \infty$. For these,
    $\braket{v}{\op{M}v} = \sum_j j\, \abs{v_j}^2$ is finite; for the vector with $v_j = 1/j$
    exactly this sum diverges, as at the end of Foundation~\ref{c03.f3}.
  \end{example}

  For an unbounded operator the matrix elements are no longer enough either. For a bounded
  operator $\op{T}$ defined on all of $\elltwo$, the numbers
  $\entry{\op{T}}{j}{k} = \vb{u}_j\adj (\op{T}\vb{u}_k)$ determine the operator. With the partial
  sums $\vb{v}_N = \sum_{k \leq N} v_k \vb{u}_k$ we have
  $\norm{\op{T}\vb{v} - \op{T}\vb{v}_N} \leq C\, \norm{\vb{v} - \vb{v}_N}$, and this goes to
  zero (Example~\ref{c03.ex.elltwo-basis}); by the Schwarz inequality the entries then converge
  too, and
  \begin{equation}\label{c03.eq.elemente-unendlich}\claim{c03.cl.elemente-unendlich}{derived}
    (\op{T}\vb{v})_j = \sum_{k=1}^{\infty} \entry{\op{T}}{j}{k}\, v_k .
  \end{equation}
  For $\op{M}$, by contrast, $\entry{\op{M}}{j}{k} = j\, \kron{j}{k}$, whether one applies
  $\op{M}$ only to the vectors with finitely many nonzero entries or to all those with
  $\sum_j j^2 \abs{v_j}^2 < \infty$. These are two different operators with the same matrix
  elements; only the domain distinguishes them.

  Such operators are no curiosity. Reed and Simon begin their chapter on unbounded operators
  with the observation that many of the most important operators of mathematical physics are
  not bounded \srcref{reed-simon-1}{p.~249}. That unbounded symmetric operators cannot be
  defined on the whole space is no accident:

  \begin{theorem}[Hellinger and Toeplitz]\label{c03.thm.hellinger-toeplitz}\claim[proof-omitted]{c03.cl.hellinger-toeplitz}{source-quoted}
    Let $\op{T}$ be an operator defined on a whole Hilbert space $\Hil$, with
    $\braket{\phi}{\op{T}\psi} = \braket{\op{T}\phi}{\psi}$ for all
    $\ket{\phi}, \ket{\psi} \in \Hil$. Then $\op{T}$ is bounded.
  \end{theorem}

  Reed and Simon prove the theorem with the closed graph theorem. They add that in quantum
  mechanics there are unbounded operators, such as the energy, that we want to satisfy this
  equation in some sense; by the theorem they cannot be defined on the whole space
  \srcref{reed-simon-1}{p.~84}. Such an operator is defined only on a subspace, its
  \emph{domain} $\dom{\op{T}}$, and then one must also watch the domains for the adjoint
  operator (\defref{D19}). Reed and Simon call a densely defined operator \emph{symmetric} (or
  Hermitian) and \emph{self-adjoint} if
  \begin{equation}\label{c03.eq.symmetrisch}\claim{c03.cl.symmetrisch}{source-quoted}
    \op{T} \subset \op{T}\adj \quad \text{\tpcifenglish{(symmetric)}{(symmetrisch)}} ,
    \qquad
    \op{T} = \op{T}\adj \quad \text{\tpcifenglish{(self-adjoint)}{(selbstadjungiert)}} .
  \end{equation}
  Here $\op{T} \subset \op{T}\adj$ means: $\dom{\op{T}} \subset \dom{\op{T}\adj}$, and on
  $\dom{\op{T}}$ the two act alike; equivalently,
  $\braket{\op{T}\phi}{\psi} = \braket{\phi}{\op{T}\psi}$ for all $\ket{\phi}, \ket{\psi}$ in
  $\dom{\op{T}}$. Self-adjoint additionally demands $\dom{\op{T}} = \dom{\op{T}\adj}$
  \srcref{reed-simon-1}{p.~255}. In this usage, Hermitian thus means the same as symmetric, and
  that is weaker than self-adjoint. The difference arises only through the restricted domain:
  for an operator defined on the whole space the notions coincide
  (Proposition~\ref{c03.prop.symmetrisch}). The difference is very important. Among the
  symmetric operators the spectral theorem holds only for the self-adjoint ones, and only
  self-adjoint operators generate the one-parameter unitary groups that give the dynamics of
  quantum mechanics \srcref{reed-simon-1}{p.~256}.

  \begin{proposition}[On the whole space, symmetric is self-adjoint]\label{c03.prop.symmetrisch}\claim[proof-included]{c03.cl.ganz}{derived}
    If $\op{T}$ is defined on all of $\Hil$ and symmetric, then $\op{T}$ is self-adjoint. In
    particular, every Hermitian operator on a finite-dimensional Hilbert space is self-adjoint.
  \end{proposition}
  \begin{proof}
    From $\dom{\op{T}} = \Hil$ and $\dom{\op{T}} \subset \dom{\op{T}\adj} \subset \Hil$ it
    follows that $\dom{\op{T}\adj} = \dom{\op{T}}$. On a finite-dimensional space every
    operator is defined on all of $\Hil$, and a Hermitian operator satisfies
    $\braket{\op{A}\phi}{\psi} = \braket{\phi}{\op{A}\psi}$ for all vectors (\defref{D18}), so it
    is symmetric.
  \end{proof}
\end{vertiefung}

Some classes are subclasses of others. The definitions claim some of this; here are the proofs.
They use only the definition of the adjoint operator and the positivity of the inner product.

\begin{proposition}[Hermitian and positive operators]\label{c03.prop.positiv}\claim[proof-included]{c03.cl.positiv}{derived}
  Let $\op{A}$ and $\op{B}$ be operators on a finite-dimensional Hilbert space $\Hil$.
  \begin{enumerate}
    \item $\op{A}$ is Hermitian if and only if $\braket{\psi}{\op{A}\psi}$ is a real number for
      every $\ket{\psi}$.
    \item Every positive operator is Hermitian.
    \item $\op{B}\adj\op{B}$ is positive.
    \item Every projector is positive.
  \end{enumerate}
\end{proposition}
\begin{proof}
  Conjugating the defining equation of \defref{D17} gives, with
  Proposition~\ref{c03.prop.symmetrie}, $\braket{\psi}{\op{A}\adj\phi} = \braket{\op{A}\psi}{\phi}$.
  On~1: if $\op{A}$ is Hermitian, then
  $\braket{\psi}{\op{A}\psi} = \braket{\op{A}\psi}{\psi} = \braket{\psi}{\op{A}\psi}^*$, so it
  is real. Conversely, let $\braket{\psi}{\op{A}\psi}$ be real for every $\ket{\psi}$, and let
  $\op{D} = \op{A} - \op{A}\adj$. Then
  $\braket{\psi}{\op{D}\psi} = \braket{\psi}{\op{A}\psi} - \braket{\op{A}\psi}{\psi} = \braket{\psi}{\op{A}\psi} - \braket{\psi}{\op{A}\psi}^* = 0$
  for every $\ket{\psi}$. Inserting $\ket{\phi} + \ket{\xi}$ and $\ket{\phi} + i\ket{\xi}$,
  the terms with $\ket{\phi}$ alone and $\ket{\xi}$ alone drop out, and what remains is
  \begin{equation}\label{c03.eq.polarisation}
    \braket{\phi}{\op{D}\xi} + \braket{\xi}{\op{D}\phi} = 0 ,
    \qquad
    i\, \braket{\phi}{\op{D}\xi} - i\, \braket{\xi}{\op{D}\phi} = 0 .
  \end{equation}
  Dividing the second equation by $i$ and adding it to the first gives
  $\braket{\phi}{\op{D}\xi} = 0$ for all $\ket{\phi}, \ket{\xi}$. With
  $\ket{\phi} = \op{D}\ket{\xi}$ and positivity, $\op{D}\ket{\xi} = 0$, so
  $\op{A}\adj = \op{A}$. On~2: for positive $\op{A}$, $\braket{\psi}{\op{A}\psi} \geq 0$, in
  particular real; continue with~1. On~3:
  $\braket{\psi}{\op{B}\adj\op{B}\psi} = \braket{\op{B}\psi}{\op{B}\psi} \geq 0$ by the
  positivity of the inner product. On~4: for a projector,
  $\op{P} = \op{P}^2 = \op{P}\adj\op{P}$; continue with~3.
\end{proof}

Part~3 is very helpful. Its proof is the same step as in the proof of the Schwarz inequality,
Equation~\eqref{c03.eq.schwarz-beweis}: the inner product of a vector with itself is never
negative. With this step the proof of Heisenberg's uncertainty relation begins in Chapter~6.
There $\op{A}$ and $\op{B}$ are Hermitian and $\lambda$ is real, and the operator
$(\op{A} + i\lambda\op{B})\adj(\op{A} + i\lambda\op{B})$ is positive because it has the form
$\op{D}\adj\op{D}$ \srcref{tpc-script-v0.2}{vl09, proof of the uncertainty relation}.

\begin{proposition}[Isometries and normal operators]\label{c03.prop.isometrie}\claim[proof-included]{c03.cl.isometrie}{derived}
  Let $\Hil$ be finite-dimensional.
  \begin{enumerate}
    \item An operator $\op{V}$ is isometric if and only if
      $\braket{\op{V}\phi}{\op{V}\psi} = \braket{\phi}{\psi}$ for all
      $\ket{\phi}, \ket{\psi} \in \Hil$.
    \item Every isometry is unitary.
    \item Hermitian and unitary operators are normal, and so are positive operators, projectors
      and isometries.
  \end{enumerate}
\end{proposition}
\begin{proof}
  On~1: $\braket{\op{V}\phi}{\op{V}\psi} = \braket{\phi}{\op{V}\adj\op{V}\psi}$. If
  $\op{V}\adj\op{V} = \Id$, this is $\braket{\phi}{\psi}$. If conversely
  $\braket{\phi}{(\op{V}\adj\op{V} - \Id)\psi} = 0$ for all $\ket{\phi}$, then with
  $\ket{\phi} = (\op{V}\adj\op{V} - \Id)\ket{\psi}$ and positivity
  $\op{V}\adj\op{V}\ket{\psi} = \ket{\psi}$. On~2: from $\op{V}\ket{\psi} = 0$ it follows by~1
  that $\norm{\psi} = \norm{\op{V}\psi} = 0$; so $\op{V}$ maps only the zero vector to zero. By
  the rank-nullity theorem of linear algebra such a linear map of a finite-dimensional space
  into itself is also surjective: every vector is of the form $\op{V}\ket{\psi}$. For it,
  $\op{V}\op{V}\adj(\op{V}\ket{\psi}) = \op{V}(\op{V}\adj\op{V})\ket{\psi} = \op{V}\ket{\psi}$,
  so $\op{V}\op{V}\adj = \Id$. On~3: for Hermitian $\op{A}$,
  $\op{A}\adj\op{A} = \op{A}\op{A} = \op{A}\op{A}\adj$; for unitary, both sides are $\Id$.
  Positive operators and projectors are Hermitian by Proposition~\ref{c03.prop.positiv},
  isometries are unitary by~2.
\end{proof}

Part~2 really needs finite dimension: in infinitely many dimensions there are isometries that
are not unitary\zoomref{c03.v6}.

\begin{vertiefung}[id=c03.v6,
    titel={Isometric but not unitary},
    wozu={In finite dimension every isometry is unitary
      (Proposition~\ref{c03.prop.isometrie}). Here is an isometry of $\elltwo$ that is not, and
      one sees at which step of the proof it fails.},
    voraussetzung={Proposition~\ref{c03.prop.isometrie}; Deeper blocks~\ref{c03.v1} and~\ref{c03.v4}
      (the example of a basis of $\elltwo$). The core up to Chapter~4 does not need this block.},
    zurueck={In infinitely many dimensions an operator can preserve the inner product without
      being unitary. Part~2 of Proposition~\ref{c03.prop.isometrie} holds only in finite
      dimension.}]
  \begin{example}[The right shift]\label{c03.ex.shift}\claim{c03.cl.shift}{derived}
    The right shift $\op{S}$ moves every entry of a vector in $\elltwo$ one place to the right
    and puts a zero in front:
    \begin{equation}\label{c03.eq.shift}
      \op{S} \begin{pmatrix} v_1 \\ v_2 \\ v_3 \\ \vdots \end{pmatrix}
      = \begin{pmatrix} 0 \\ v_1 \\ v_2 \\ \vdots \end{pmatrix} ,
      \qquad
      \op{S}\adj \begin{pmatrix} w_1 \\ w_2 \\ w_3 \\ \vdots \end{pmatrix}
      = \begin{pmatrix} w_2 \\ w_3 \\ w_4 \\ \vdots \end{pmatrix} .
    \end{equation}
    The series $(\op{S}\vb{v})\adj(\op{S}\vb{w})$ has the same terms as $\vb{v}\adj\vb{w}$, only
    one place later; so $\op{S}\vb{v} \in \elltwo$, and $\op{S}$ preserves the inner product.
    The adjoint operator is read off from
    $\vb{w}\adj(\op{S}\vb{v}) = \sum_{k \geq 1} w_{k+1}^*\, v_k$: it shifts to the left and
    drops the first entry. Hence $\op{S}\adj\op{S} = \Id$, but
    $\op{S}\op{S}\adj \vb{u}_1 = 0$ for the first basis vector $\vb{u}_1$, so
    $\op{S}\op{S}\adj \neq \Id$. $\op{S}$ is isometric, but neither unitary nor normal.
  \end{example}

  It fails at the second step of the proof of Proposition~\ref{c03.prop.isometrie}: $\op{S}$
  maps only the zero vector to zero, but it is not surjective, because every vector
  $\op{S}\vb{v}$ has a zero as its first entry. Reed and Simon use the same example: the right
  shift on $\elltwo$ is not unitary, because the vector $(1, 0, 0, \dots)$ is not in its range
  \srcref{reed-simon-1}{p.~197}.
\end{vertiefung}

\section{Matrix elements and ket-bras}\label{c03.s11}

With the help of a basis one can describe operators by numbers. Let
$\{\ket{\phi_x}\}_{x \in X}$ be a basis of $\Hil$ with a finite index set $X$.

\tpcdef{D93}

Where the formula for $\op{A}\ket{\Psi}$ in \defref{D93} comes from, we work out once; this is
the exercise we do every day in quantum mechanics. We first expand $\ket{\Psi}$ in the basis,
move the numbers $\braket{\phi_y}{\Psi}$ past $\op{A}$, since numbers always commute with
operators, and then expand each vector $\op{A}\ket{\phi_y}$ once more:
\begin{equation}\label{c03.eq.wirkung}\claim{c03.cl.wirkung}{derived}
  \begin{aligned}
    \op{A}\ket{\Psi}
    &= \op{A} \sum_{y \in X} \braket{\phi_y}{\Psi}\, \ket{\phi_y}
     = \sum_{y \in X} \braket{\phi_y}{\Psi}\, \op{A}\ket{\phi_y} \\
    &= \sum_{y \in X} \braket{\phi_y}{\Psi} \sum_{x \in X} \braket{\phi_x}{\op{A}\phi_y}\, \ket{\phi_x}
     = \sum_{x, y \in X} \ket{\phi_x}\, \entry{\op{A}}{x}{y}\, \braket{\phi_y}{\Psi} .
  \end{aligned}
\end{equation}
The sums over $x$ and $y$ run independently of each other.

Note that nowhere in this equation does a matrix occur, only basis vectors and numbers. If one
knows all matrix elements, one can apply $\op{A}$ to every vector: an operator is known as soon
as one knows its matrix elements. In the finite-dimensional case one simply writes them into a
matrix. In infinitely many dimensions there is no finite matrix any more, but there are the
numbers $\entry{\op{A}}{x}{y}$, a table with infinitely many rows and columns. There the sums
become series, and that the matrix elements determine the operator holds only under
hypotheses: for bounded operators defined on the whole space, but not without more for
unbounded ones\zoomref{c03.v5}.

The adjoint operator, too, can be read off from the matrix elements. With the conjugated
equation from the proof of Proposition~\ref{c03.prop.positiv},
\begin{equation}\label{c03.eq.adj-elemente}\claim{c03.cl.adj-elemente}{derived}
  \entry{\op{A}\adj}{x}{y}
  = \braket{\phi_x}{\op{A}\adj\phi_y}
  = \braket{\op{A}\phi_x}{\phi_y}
  = \braket{\phi_y}{\op{A}\phi_x}^*
  = \entry{\op{A}}{y}{x}^* .
\end{equation}
Forming the adjoint is thus, in every orthonormal basis, transposition with complex
conjugation, for every inner product. Reed and Simon remind the reader that a linear
transformation on $\mathbb{C}^n$ is self-adjoint if and only if its matrix in an orthonormal
basis is unchanged by reflection across the diagonal followed by conjugation
\srcref{reed-simon-1}{p.~187}.

\subsection{Kets and bras as operators}\label{c03.s11.1}

So far the ket $\ket{\Psi}$ simply stood for a vector $\vb{\Psi} \in \Hil$. It pays to make a
very subtle distinction here, because it simplifies what follows. From now on the ket is the
linear map
\begin{equation}\label{c03.eq.ket-op}
  \ket{\Psi}\colon \mathbb{C} \to \Hil ,
  \qquad
  z \mapsto z\, \vb{\Psi} .
\end{equation}
So a ket eats a complex number and produces a vector. I have always left the number out so far;
what was meant was $z = 1$.

The equation of \defref{D17} defines the adjoint operator word for word also for a linear map
from one Hilbert space into another, here from $\mathbb{C}$ to $\Hil$; here
$\mathbb{C} = \mathbb{C}^1$ is a Hilbert space with the inner product $a^* b$. One checks that
the adjoint of the ket is the bra:
\begin{equation}\label{c03.eq.bra-op}\claim{c03.cl.bra-op}{derived}
  \bra{\Psi} = (\ket{\Psi})\adj \colon \Hil \to \mathbb{C} ,
  \qquad
  \vb{\Phi} \mapsto \braket{\Psi}{\Phi} .
\end{equation}
For every number $z$ and every vector $\vb{\Phi}$ we have
$\braket{\Phi}{z\Psi} = z\, \braket{\Phi}{\Psi} = \braket{\Psi}{\Phi}^* z$, and on the right
stands the inner product of the number $\braket{\Psi}{\Phi}$ with the number $z$ in
$\mathbb{C}$. So a bra eats a vector and produces a complex number, exactly the other way round
from a ket. This fits Section~\ref{c03.s4}, where the bra was the assignment that gives every
ket a number.

If one multiplies a bra by a ket, one composes two maps: the ket turns a number into a vector,
the bra turns it back into a number. The result is the map from $\mathbb{C}$ to $\mathbb{C}$
\begin{equation*}
  z \mapsto \bra{\Phi}\, \ket{\Psi}\, z = \braket{\Phi}{\Psi}\, z .
\end{equation*}
If we make no distinction in the notation between a linear map from $\mathbb{C}$ to
$\mathbb{C}$ and the number it multiplies by, then $\bra{\Phi}\, \ket{\Psi}$ and
$\braket{\Phi}{\Psi}$ mean the same. That is the main message of this section. The
justification was tedious; you can now forget it and use this identification from now on. For
calculating nothing changes; the distinction matters only for the interpretation.

In the other order, ket times bra, an operator on $\Hil$ arises: the bra eats a vector and
returns a number, the ket eats this number and returns a vector. Such operators are extremely
useful in quantum mechanics.

\tpcdef{D94}

We go one step further in simplifying. As soon as a basis $\{\ket{\phi_x}\}_{x \in X}$ is fixed
and clear from the context, the letter $\phi$ carries no information any more; all the
information is in the index. We drop it and write $\ket{x}$ and $\bra{x}$. With this the
expansion formulas for vectors and for operators become operator equations.

\begin{proposition}[Calculating with ket-bras]\label{c03.prop.ketbra}\claim[proof-included]{c03.cl.ketbra}{derived}
  For vectors $\ket{\Psi}, \ket{\Phi}$, a basis $\{\ket{x}\}_{x \in X}$ with a finite index set
  and every operator $\op{A}$,
  \begin{equation}\label{c03.eq.ketbra}
    (\ketbra{\Psi}{\Phi})\adj = \ketbra{\Phi}{\Psi} ,
    \qquad
    \Id = \sum_{x \in X} \ketbra{x}{x} ,
    \qquad
    \op{A} = \sum_{x, y \in X} \melem{x}{\op{A}}{y}\, \ketbra{x}{y} .
  \end{equation}
\end{proposition}
\begin{proof}
  We apply both sides to an arbitrary vector; that is the general pattern for operator
  equations. First equation: for all $\ket{\zeta}, \ket{\xi}$,
  $\braket{\zeta}{(\ketbra{\Psi}{\Phi})\xi} = \braket{\Phi}{\xi}\, \braket{\zeta}{\Psi}$, and
  $\braket{(\ketbra{\Phi}{\Psi})\zeta}{\xi} = \braket{\Psi}{\zeta}^*\, \braket{\Phi}{\xi}$ is the
  same; by Proposition~\ref{c03.prop.adjungiert}, $\ketbra{\Phi}{\Psi}$ is the adjoint
  operator. Second equation: on the left stands $\Id\ket{\Psi} = \ket{\Psi}$, on the right by
  \defref{D94} $\sum_x \braket{x}{\Psi}\, \ket{x}$, and that is $\ket{\Psi}$ by
  Proposition~\ref{c03.prop.entwicklung}. Third equation: on the right stands
  $\sum_{x,y} \melem{x}{\op{A}}{y}\, \braket{y}{\Psi}\, \ket{x}$, and that is
  $\op{A}\ket{\Psi}$ by Equation~\eqref{c03.eq.wirkung}.
\end{proof}

The equation $\Id = \sum_x \ketbra{x}{x}$ is one of the most helpful equations of quantum
mechanics; I still use it regularly today. Physically it contains no information: it only says
that one equals one. Mathematically, however, it makes life enormously easier, because one may
insert this complicated one at any point of a calculation. Matrix elements are now simply
written $\melem{x}{\op{A}}{y}$.

\section{Three operators on $\mathbb{C}^3$}\label{c03.s12}

Finally, an example in which everything comes together.

\begin{example}[Three operators on $\mathbb{C}^3$]\label{c03.ex.drei}\claim{c03.cl.drei}{derived}
  Let $\Hil = \mathbb{C}^3$ with the inner product of Section~\ref{c03.s2}, and let
  \begin{equation}\label{c03.eq.abc}
    \op{A} = \begin{pmatrix} 0 & 1 & 0 \\ 1 & 0 & 1 \\ 0 & 1 & 0 \end{pmatrix} ,
    \qquad
    \op{B} = \frac13 \begin{pmatrix} 1 & 1 & 1 \\ 1 & 1 & 1 \\ 1 & 1 & 1 \end{pmatrix} ,
    \qquad
    \op{C} = \begin{pmatrix} 0 & 1 & 0 \\ 0 & 0 & i \\ -1 & 0 & 0 \end{pmatrix} .
  \end{equation}
  Transposing and conjugating gives the same matrix for $\op{A}$ and $\op{B}$; both are
  Hermitian, $\op{A}\adj = \op{A}$ and $\op{B}\adj = \op{B}$. $\op{B}$ is also a projector; that
  takes a little more work, but not much. Every entry of $\op{B}^2$ is
  $\tfrac13 \cdot \tfrac13 \, (1 \cdot 1 + 1 \cdot 1 + 1 \cdot 1) = \tfrac13$, so
  $\op{B}^2 = \op{B}$, and by Proposition~\ref{c03.prop.positiv} $\op{B}$ is positive.
  $\op{A}$, on the other hand, is \emph{not} positive: for the vector $\ket{\psi}$ with the
  entries $\tfrac12$, $-\tfrac{1}{\sqrt{2}}$, $\tfrac12$,
  \begin{equation}\label{c03.eq.nicht-positiv}
    \melem{\psi}{\op{A}}{\psi}
    = \begin{pmatrix} \frac12 & -\frac{1}{\sqrt{2}} & \frac12 \end{pmatrix}
      \begin{pmatrix} 0 & 1 & 0 \\ 1 & 0 & 1 \\ 0 & 1 & 0 \end{pmatrix}
      \begin{pmatrix} \frac12 \\ -\frac{1}{\sqrt{2}} \\ \frac12 \end{pmatrix}
    = \begin{pmatrix} \frac12 & -\frac{1}{\sqrt{2}} & \frac12 \end{pmatrix}
      \begin{pmatrix} -\frac{1}{\sqrt{2}} \\ 1 \\ -\frac{1}{\sqrt{2}} \end{pmatrix}
    = -\sqrt{2} .
  \end{equation}
  $\op{C}$ is unitary, because
  \begin{equation}\label{c03.eq.c-unitaer}
    \op{C}\adj\op{C}
    = \begin{pmatrix} 0 & 0 & -1 \\ 1 & 0 & 0 \\ 0 & -i & 0 \end{pmatrix}
      \begin{pmatrix} 0 & 1 & 0 \\ 0 & 0 & i \\ -1 & 0 & 0 \end{pmatrix}
    = \begin{pmatrix} 1 & 0 & 0 \\ 0 & 1 & 0 \\ 0 & 0 & 1 \end{pmatrix} ,
  \end{equation}
  and likewise $\op{C}\op{C}\adj = \Id$. With the canonical basis $\ket{0}, \ket{1}, \ket{2}$
  of $\mathbb{C}^3$ (the columns with the $1$ in the first, second, third entry) all nine
  matrix elements of $\op{B}$ are equal, $\melem{x}{\op{B}}{y} = \tfrac13$, and
  \begin{equation}\label{c03.eq.c-ketbra}
    \op{C} = \ketbra{0}{1} + i\, \ketbra{1}{2} - \ketbra{2}{0} .
  \end{equation}
  One checks this by applying both sides to the basis vectors. The matrix maps $\ket{1}$ to
  $\ket{0}$; on the right, of the three terms only $\braket{1}{1}\, \ket{0} = \ket{0}$ remains,
  because $\braket{2}{1} = \braket{0}{1} = 0$. Likewise $\ket{0}$ is mapped to $-\ket{2}$ and
  $\ket{2}$ to $i\ket{1}$. The order within each term matters.
\end{example}

So one can do everything with matrices or everything with ket-bras. In the following chapters I
mix both representations; in the end you should master both. Which one is best suited usually
depends on the dimension: for finite-dimensional systems one likes to calculate with matrices,
for infinite-dimensional ones almost only with the ket notation.

\begin{exercise}[Unitary, with ket-bras]\label{c03.exc.ketbra}\claim{c03.cl.c-ketbra}{derived}
  Compute $\op{C}\adj$ from Equation~\eqref{c03.eq.c-ketbra} with
  Proposition~\ref{c03.prop.ketbra}, and show $\op{C}\adj\op{C} = \Id$ without writing down a
  matrix.
\end{exercise}
\begin{solution}
  By the rules of \defref{D17} and Proposition~\ref{c03.prop.ketbra},
  $\op{C}\adj = \ketbra{1}{0} - i\, \ketbra{2}{1} - \ketbra{0}{2}$. In the product
  $\op{C}\adj\op{C}$, $\ketbra{a}{b}\, \ketbra{c}{d} = \braket{b}{c}\, \ketbra{a}{d}$, and this is
  zero for $b \neq c$. Of the nine products, three remain:
  \begin{equation}\label{c03.eq.c-ketbra-unitaer}
    \op{C}\adj\op{C}
    = \ketbra{1}{1} + (-i)\, i\, \ketbra{2}{2} + (-1)(-1)\, \ketbra{0}{0}
    = \ketbra{0}{0} + \ketbra{1}{1} + \ketbra{2}{2}
    = \Id .
  \end{equation}
  The last equation is the completeness relation of Proposition~\ref{c03.prop.ketbra}.
\end{solution}

What the three calculations have in common becomes visible when one plots the eigenvalues. An
eigenvalue is, informally, a number $\lambda$ by which $\op{A}$ simply multiplies a vector
$\ket{v} \neq 0$, $\op{A}\ket{v} = \lambda\ket{v}$; the precise definition, and how one finds
all eigenvalues, comes later in this chapter. For every eigenvalue the defining equation of the
class carries over to the number $\lambda$:
\begin{equation}\label{c03.eq.orte}\claim{c03.cl.orte}{derived}
  \begin{aligned}
    &\op{A}\adj = \op{A} &&\Rightarrow\quad \lambda^* = \lambda , \\
    &\text{\tpcifenglish{for all}{für alle} } \ket{\psi}\colon \braket{\psi}{\op{A}\psi} \geq 0 &&\Rightarrow\quad \lambda \geq 0 , \\
    &\op{A}^2 = \op{A}\adj = \op{A} &&\Rightarrow\quad \lambda^2 = \lambda , \\
    &\op{A}\adj\op{A} = \Id &&\Rightarrow\quad \abs{\lambda} = 1 .
  \end{aligned}
\end{equation}
One may assume $\norm{v} = 1$. Then $\melem{v}{\op{A}}{v} = \lambda$; this is real for
Hermitian $\op{A}$ (Proposition~\ref{c03.prop.positiv}) and non-negative for positive $\op{A}$.
For a projector, $\lambda^2\ket{v} = \op{A}^2\ket{v} = \op{A}\ket{v} = \lambda\ket{v}$, and for
an isometry $\abs{\lambda} = \norm{\op{A}v} = \norm{v} = 1$. So the Hermitian operators have
their eigenvalues on the real axis, the positive ones on the half-axis $\lambda \geq 0$, the
projectors in the two points $0$ and $1$, the unitary ones on the unit circle. Normality alone
does not constrain the eigenvalues: for every complex number $\lambda$, $\lambda\Id$ is normal
and has the eigenvalue $\lambda$. For the three operators of the example one finds eigenvectors
by inserting; for $\op{A}$, for instance, the vectors with entries $1, \sqrt{2}, 1$ and
$1, 0, -1$ and $1, -\sqrt{2}, 1$, for $\op{B}$ the vectors $1, 1, 1$ and $1, -1, 0$, for
$\op{C}$ the vectors $1, \lambda, -i\lambda^2$. The eigenvalues are
\begin{equation}\label{c03.eq.spektren}\claim{c03.cl.spektren}{derived}
  \begin{aligned}
    &\op{A}\colon && 0 ,\ \sqrt{2} ,\ -\sqrt{2} , \\
    &\op{B}\colon && 1 ,\ 0 , \\
    &\op{C}\colon && e^{-i\pi/6} ,\ i ,\ e^{7i\pi/6} , \qquad \text{\tpcifenglish{the three solutions of}{die drei Lösungen von}}\ \lambda^3 = -i ,
  \end{aligned}
\end{equation}
and there are no others. Figure~\ref{c03.fig.spektren} records in part (a) the places where the
eigenvalues of a class may lie, and in part (b) enters the eigenvalues of the three example
operators. In (b), $-\sqrt{2}$ lies to the left of zero; precisely there the positivity of
$\op{A}$ fails. The vector $\ket{\psi}$ of Equation~\eqref{c03.eq.nicht-positiv} is the
normalised eigenvector for this eigenvalue, and that is why
$\melem{\psi}{\op{A}}{\psi} = -\sqrt{2}$.

\begin{tpcfigure}{c03.fig.spektren}
  \caption{Operator classes recognised by their eigenvalues. (a) Where in the complex plane the
    eigenvalues of a class may lie (Equation~\eqref{c03.eq.orte}): Hermitian on the real axis
    (solid line), positive on the half-axis $\lambda \geq 0$ (broad strip), projector in the two
    points $0$ and $1$ (filled dots), unitary on the unit circle (dashed). Normality alone admits
    every complex number. The nesting projector, positive, Hermitian is visible: two points on a
    half-axis on a line. (b) The eigenvalues of the operators of Example~\ref{c03.ex.drei}
    (Equation~\eqref{c03.eq.spektren}): $\op{A}$ (filled discs) at $0$ and $\pm\sqrt{2}$,
    $\op{B}$ (rings) at $0$ and $1$, $\op{C}$ (diamonds) at three points on the unit circle that
    divide it into three equal arcs. The eigenvalue $-\sqrt{2}$ of $\op{A}$ lies to the left of
    zero: $\op{A}$ is Hermitian but not positive.}
\end{tpcfigure}

% ── Fundament ─────────────────────────────────────────────────────────────────────────
% New sections of the chapter (vl05) go BEFORE this line.

\begin{fundament}[id=c03.f1,
    titel={The Cauchy-Schwarz inequality},
    wozu={Block~\ref{c03.v1} shows with a rough estimate that the inner product on $\elltwo$
      is finite. Here is the sharp bound: the modulus of the inner product is at most the
      product of the lengths.},
    voraussetzung={Proposition~\ref{c03.prop.skalarprodukt}, the Schwarz inequality
      (Proposition~\ref{c03.prop.schwarz}) and block~\ref{c03.v1}. Neither the core nor the
      deeper blocks need this block.},
    zurueck={For two vectors, $\abs{\vb{v}\adj\vb{w}}^2$ is never larger than
      $(\vb{v}\adj\vb{v})(\vb{w}\adj\vb{w})$, in $\mathbb{C}^n$ as in $\elltwo$.}]
  \begin{lemma}[Cauchy-Schwarz in $\mathbb{C}^n$]\label{c03.lem.cs}\claim[proof-linked]{c03.cl.cs}{derived}
    For all $\vb{v}, \vb{w} \in \mathbb{C}^n$,
    \begin{equation}\label{c03.eq.cs}
      \abs{\vb{v}\adj\vb{w}}^2 \leq (\vb{v}\adj\vb{v})\,(\vb{w}\adj\vb{w}) .
    \end{equation}
  \end{lemma}

  The lemma is the Schwarz inequality of Proposition~\ref{c03.prop.schwarz} for the inner
  product $\vb{v}\adj\vb{w}$ of $\mathbb{C}^n$. This inner product satisfies the conditions of
  \defref{D82}, since these are rules~2 to~4 of Proposition~\ref{c03.prop.skalarprodukt}; with
  $\norm{v}^2 = \vb{v}\adj\vb{v}$, Equation~\eqref{c03.eq.cs} is the square of
  Equation~\eqref{c03.eq.schwarz}. Reed and Simon prove the same inequality for every vector
  space with an inner product, by a different route \srcref{reed-simon-1}{p.~38}. For
  $\elltwo$ it follows here from the lemma by a limit.

  \begin{proposition}[Cauchy-Schwarz in $\elltwo$]\label{c03.prop.cs-elltwo}\claim[proof-included]{c03.cl.cs-elltwo}{derived}
    For all $\vb{v}, \vb{w} \in \elltwo$,
    $\abs{\vb{v}\adj\vb{w}}^2 \leq (\vb{v}\adj\vb{v})\,(\vb{w}\adj\vb{w})$.
  \end{proposition}
  \begin{proof}
    For every $N$, Lemma~\ref{c03.lem.cs} applied to the first $N$ entries is the inequality
    $\abs{\sum_{j \leq N} v_j^* w_j}^2 \leq \sum_{j \leq N} \abs{v_j}^2 \sum_{j \leq N} \abs{w_j}^2$.
    The right-hand side is at most $(\vb{v}\adj\vb{v})(\vb{w}\adj\vb{w})$. The left-hand side
    tends to $\abs{\vb{v}\adj\vb{w}}^2$ as $N \to \infty$, because the series converges by
    Proposition~\ref{c03.prop.elltwo}.
  \end{proof}
\end{fundament}

\begin{fundament}[id=c03.f2,
    titel={Functions that differ only on a null set},
    wozu={Block~\ref{c03.v2} leaves open whether $\braket{f}{f} = 0$ holds only for $f = 0$.
      For functions it does not. Here is the repair and what it means for the value of a
      function at a single point.},
    voraussetzung={Block~\ref{c03.v2}. Lebesgue measure and the Lebesgue integral are used by
      name only.},
    zurueck={Two functions that differ only at single points are the same vector in
      $\Ltwo{\mathbb{R}}$. The value at a single point has no meaning, but integrals over
      intervals do.}]
  \begin{example}[A function of length zero]\label{c03.ex.punkt}\claim{c03.cl.punkt}{derived}
    Let $f(0) = 1$ and $f(x) = 0$ for $x \neq 0$. This function is piecewise continuous and
    not the zero function, but
    \begin{equation}\label{c03.eq.punkt}
      \braket{f}{f} = \int_{-\infty}^{\infty} \abs{f(x)}^2\, \dd x = 0 ,
    \end{equation}
    because an integrand that is nonzero at a single point only has integral zero. Rule~4 of
    Proposition~\ref{c03.prop.skalarprodukt} therefore fails for functions.
  \end{example}

  The repair is to identify such functions with the zero function. In the precise version,
  $\Ltwo{\mathbb{R}}$ is therefore a space of classes of functions.

  \tpcdef{D80}

  For square-integrable $f$ and $g$,
  \begin{equation}\label{c03.eq.klassen}\claim{c03.cl.klassen}{derived}
    f \sim g
    \quad \Leftrightarrow \quad
    \int_{-\infty}^{\infty} \abs{f(x) - g(x)}^2\, \dd x = 0 .
  \end{equation}
  Reed and Simon say for integrable functions that $f(x) = g(x)$ almost everywhere is the same
  as $\int \abs{f - g}\, \dd x = 0$ \srcref{reed-simon-1}{p.~17}. Applied to the integrable
  function $\abs{f - g}^2$ and the zero function, this gives
  Equation~\eqref{c03.eq.klassen}, since $\abs{f - g}^2$ is zero exactly where $f = g$. It
  also follows that the inner product does not depend on which function one takes from a
  class: if $n \sim 0$, then $\abs{n^* g}$ is zero almost everywhere, so by the same statement
  $\int \abs{n^* g}\, \dd x = 0$ and hence $\braket{n}{g} = 0$. In $\Ltwo{\mathbb{R}}$ rule~4
  now holds as well: $\braket{f}{f} = 0$ implies $f \sim 0$, and that is the zero vector.

  The price is that an element of $\Ltwo{\mathbb{R}}$ has no value at a single point: if one
  changes a function at one point, it stays in its class. Reed and Simon put it like this for
  integrable functions: for an element $f$, the symbol $f(x)$ at a particular point $x$ does
  not make sense \srcref{reed-simon-1}{p.~17}. For functions that have a continuous
  representative, the following statement helps:

  \begin{proposition}[Continuous representatives]\label{c03.prop.stetig}\claim[proof-omitted]{c03.cl.stetig}{source-quoted}
    If a class contains a continuous function, it contains no other continuous function. One
    may then write $f(x)$ for the value of this continuous function.
  \end{proposition}

  Reed and Simon state this for the classes of integrable functions
  \srcref{reed-simon-1}{p.~18}; it depends only on the relation “equal almost everywhere”,
  which is the same for $\Ltwo{\mathbb{R}}$.
\end{fundament}

\begin{fundament}[id=c03.f3,
    titel={Completeness},
    wozu={The core says that a vector space with an inner product lacks one more condition to
      be a Hilbert space and that it always holds in $\mathbb{C}^n$. Here is what the
      condition demands, why $\mathbb{C}^n$, $\elltwo$ and $\Ltwo{\mathbb{R}}$ satisfy it, and
      a space that violates it.},
    voraussetzung={Blocks~\ref{c03.v1} and~\ref{c03.v2}, Foundations~\ref{c03.f2}; sequences
      and limits.},
    zurueck={$\mathbb{C}^n$, $\elltwo$ and $\Ltwo{\mathbb{R}}$ are complete. Whoever computes
      with them need not think about the condition.}]
  \tpcdef{D79}

  The model is the difference between the rational and the real numbers: a sequence of
  rational numbers approaching an irrational number is a Cauchy sequence but has no limit
  among the rational numbers. $\mathbb{R}$ is complete, $\mathbb{Q}$ is not
  \srcref{reed-simon-1}{p.~6}.

  \begin{proposition}[$\mathbb{C}^n$ is complete]\label{c03.prop.cn-vollstaendig}\claim[proof-included]{c03.cl.cn-vollstaendig}{derived}
    Every Cauchy sequence in $\mathbb{C}^n$ converges to a vector in $\mathbb{C}^n$.
  \end{proposition}
  \begin{proof}
    Let $\vb{v}^{(1)}, \vb{v}^{(2)}, \dots$ be a Cauchy sequence. For every index $j$,
    \begin{equation}\label{c03.eq.komponente}
      \abs{v_j^{(m)} - v_j^{(k)}}^2
      \leq \sum_{l=1}^{n} \abs{v_l^{(m)} - v_l^{(k)}}^2
      = d\bigl(\vb{v}^{(m)}, \vb{v}^{(k)}\bigr)^2 .
    \end{equation}
    So the real and the imaginary parts of the $j$-th entries are Cauchy sequences of real
    numbers; since $\mathbb{R}$ is complete, they have limits. Let $v_j$ be the combined limit
    and $\vb{v}$ the vector with these entries. Then
    $d(\vb{v}^{(m)}, \vb{v})^2 = \sum_{j=1}^{n} \abs{v_j^{(m)} - v_j}^2$ is a sum of $n$ terms,
    each of which goes to zero.
  \end{proof}

  That the condition is not empty is shown by a space that lies between $\mathbb{C}^n$ and
  $\elltwo$.

  \begin{example}[Terminating sequences are not complete]\label{c03.ex.endlich}\claim{c03.cl.endlich}{derived}
    Let $F$ be the set of vectors with infinitely many entries of which only finitely many
    are nonzero, with the inner product of $\elltwo$. Sums and multiples of elements of $F$
    lie in $F$ again. Let $\vb{v}^{(N)}$ be the vector with entries
    $1, \tfrac12, \dots, \tfrac1N$ and only zeros after them. For $M > N$, by
    Equation~\eqref{c03.eq.rate},
    \begin{equation}\label{c03.eq.endlich}
      d\bigl(\vb{v}^{(M)}, \vb{v}^{(N)}\bigr)^2
      = \sum_{j=N+1}^{M} \frac{1}{j^2}
      \leq \sum_{j=N+1}^{M} \Bigl( \frac{1}{j-1} - \frac{1}{j} \Bigr)
      = \frac{1}{N} - \frac{1}{M} < \frac{1}{N} ,
    \end{equation}
    so the sequence is a Cauchy sequence. If it had a limit $\vb{u}$ in $F$, then by the
    argument of Equation~\eqref{c03.eq.komponente}, for every $j$,
    $u_j = \lim_{N} v_j^{(N)} = 1/j$, and $\vb{u}$ would have infinitely many nonzero
    entries. So the sequence does not converge in $F$. In $\elltwo$, by contrast, it
    converges, to the vector with $v_j = 1/j$ of Example~\ref{c03.ex.rate}: the same
    computation with $M \to \infty$ gives a distance of at most $1/\sqrt{N}$.
  \end{example}

  The space $\elltwo$ thus fills exactly the gaps of $F$; Reed and Simon say more generally
  that the terminating sequences are dense in these spaces \srcref{reed-simon-1}{p.~69}.

  \begin{theorem}[Completeness of $\elltwo$ and $\Ltwo{\mathbb{R}}$]\label{c03.thm.vollstaendig}\claim[proof-omitted]{c03.cl.vollstaendig}{source-quoted}
    The space $\elltwo$ is complete. The space $\Ltwo{\mathbb{R}}$ in the version of
    \defref{D80} is complete (theorem of Riesz and Fischer).
  \end{theorem}

  The source is Reed and Simon: for $\elltwo$ as the case $p = 2$ of the sequence spaces they
  consider for every $p \geq 1$, for $\Ltwo{\mathbb{R}}$ as the case $p = 2$ of the theorem of
  Riesz and Fischer for the spaces $L^p$ \srcref{reed-simon-1}{pp.~68--69}. For
  $\Ltwo{\mathbb{R}}$ the theorem holds for the version of Foundations~\ref{c03.f2}, with
  measurable functions and the Lebesgue integral.

  We shall resist the temptation to look at limiting cases. I shall carry out only legal
  computations in this lecture, at least I hope so. Physicists handle these rules a little
  more loosely than mathematicians anyway: as soon as something becomes infinite, we say we
  have made a mistake. A prediction must not be infinite, so we drop that computation. But
  if you want to go a little in the mathematical direction, these conditions are really very
  important.

  Mathematically this is not a rule. A vector can lie in $\elltwo$ and still produce an
  infinite sum: for $v_j = 1/j$ the sum $\sum_j \abs{v_j}^2$ is finite, but
  $\sum_j j\, \abs{v_j}^2 = \sum_j 1/j$ diverges (Example~\ref{c03.ex.rate}). Sums of this
  kind appear later in expectation values. Whether such a computation is allowed must then be
  checked; the tool for that is the domains of operators.
\end{fundament}
